\documentclass{report}

%%%%%%%%%%%%%%%%%%%%%%%%%%%%%%%%%
% PACKAGE IMPORTS
%%%%%%%%%%%%%%%%%%%%%%%%%%%%%%%%%


\usepackage[tmargin=2cm,rmargin=1in,lmargin=1in,margin=0.85in,bmargin=2cm,footskip=.2in]{geometry}
\usepackage{amsmath,amsfonts,amsthm,amssymb,mathtools}
\usepackage[varbb]{newpxmath}
\usepackage{xfrac}
\usepackage[makeroom]{cancel}
\usepackage{mathtools}
\usepackage{bookmark}
\usepackage{enumitem}
\usepackage{hyperref,theoremref}
\hypersetup{
	pdftitle={Assignment},
	colorlinks=true, linkcolor=doc!90,
	bookmarksnumbered=true,
	bookmarksopen=true
}
\usepackage[most,many,breakable]{tcolorbox}
\usepackage{xcolor}
\usepackage{varwidth}
\usepackage{varwidth}
\usepackage{etoolbox}
%\usepackage{authblk}
\usepackage{nameref}
\usepackage{multicol,array}
\usepackage{tikz-cd}
\usepackage[ruled,vlined,linesnumbered]{algorithm2e}
\usepackage{comment} % enables the use of multi-line comments (\ifx \fi) 
\usepackage{import}
\usepackage{xifthen}
\usepackage{pdfpages}
\usepackage{transparent}
\usepackage[utf8]{inputenc}
\usepackage{fontawesome}


\newcommand\mycommfont[1]{\footnotesize\ttfamily\textcolor{blue}{#1}}
\SetCommentSty{mycommfont}
\newcommand{\incfig}[1]{%
    \def\svgwidth{\columnwidth}
    \import{./figures/}{#1.pdf_tex}
}

\usepackage{tikzsymbols}
\renewcommand{\qedsymbol}{\faWhatsapp}


%\usepackage{import}
%\usepackage{xifthen}
%\usepackage{pdfpages}
%\usepackage{transparent}


%%%%%%%%%%%%%%%%%%%%%%%%%%%%%%
% SELF MADE COLORS
%%%%%%%%%%%%%%%%%%%%%%%%%%%%%%



\definecolor{myg}{RGB}{56, 140, 70}
\definecolor{myb}{RGB}{45, 111, 177}
\definecolor{myr}{RGB}{199, 68, 64}
\definecolor{mytheorembg}{HTML}{F2F2F9}
\definecolor{mytheoremfr}{HTML}{00007B}
\definecolor{mylenmabg}{HTML}{FFFAF8}
\definecolor{mylenmafr}{HTML}{983b0f}
\definecolor{mypropbg}{HTML}{f2fbfc}
\definecolor{mypropfr}{HTML}{191971}
\definecolor{myexamplebg}{HTML}{F2FBF8}
\definecolor{myexamplefr}{HTML}{88D6D1}
\definecolor{myexampleti}{HTML}{2A7F7F}
\definecolor{mydefinitbg}{HTML}{E5E5FF}
\definecolor{mydefinitfr}{HTML}{3F3FA3}
\definecolor{notesgreen}{RGB}{0,162,0}
\definecolor{myp}{RGB}{197, 92, 212}
\definecolor{mygr}{HTML}{2C3338}
\definecolor{myred}{RGB}{127,0,0}
\definecolor{myyellow}{RGB}{169,121,69}
\definecolor{myexercisebg}{HTML}{F2FBF8}
\definecolor{myexercisefg}{HTML}{88D6D1}


%%%%%%%%%%%%%%%%%%%%%%%%%%%%
% TCOLORBOX SETUPS
%%%%%%%%%%%%%%%%%%%%%%%%%%%%

\setlength{\parindent}{1cm}
%================================
% THEOREM BOX
%================================

\tcbuselibrary{theorems,skins,hooks}
\newtcbtheorem[number within=section]{Theorem}{Theorem}
{%
	enhanced,
	breakable,
	colback = mytheorembg,
	frame hidden,
	boxrule = 0sp,
	borderline west = {2pt}{0pt}{mytheoremfr},
	sharp corners,
	detach title,
	before upper = \tcbtitle\par\smallskip,
	coltitle = mytheoremfr,
	fonttitle = \bfseries\sffamily,
	description font = \mdseries,
	separator sign none,
	segmentation style={solid, mytheoremfr},
}
{th}

\tcbuselibrary{theorems,skins,hooks}
\newtcbtheorem[number within=chapter]{theorem}{Theorem}
{%
	enhanced,
	breakable,
	colback = mytheorembg,
	frame hidden,
	boxrule = 0sp,
	borderline west = {2pt}{0pt}{mytheoremfr},
	sharp corners,
	detach title,
	before upper = \tcbtitle\par\smallskip,
	coltitle = mytheoremfr,
	fonttitle = \bfseries\sffamily,
	description font = \mdseries,
	separator sign none,
	segmentation style={solid, mytheoremfr},
}
{th}


\tcbuselibrary{theorems,skins,hooks}
\newtcolorbox{Theoremcon}
{%
	enhanced
	,breakable
	,colback = mytheorembg
	,frame hidden
	,boxrule = 0sp
	,borderline west = {2pt}{0pt}{mytheoremfr}
	,sharp corners
	,description font = \mdseries
	,separator sign none
}

%================================
% Corollery
%================================
\tcbuselibrary{theorems,skins,hooks}
\newtcbtheorem[number within=section]{Corollary}{Corollary}
{%
	enhanced
	,breakable
	,colback = myp!10
	,frame hidden
	,boxrule = 0sp
	,borderline west = {2pt}{0pt}{myp!85!black}
	,sharp corners
	,detach title
	,before upper = \tcbtitle\par\smallskip
	,coltitle = myp!85!black
	,fonttitle = \bfseries\sffamily
	,description font = \mdseries
	,separator sign none
	,segmentation style={solid, myp!85!black}
}
{th}
\tcbuselibrary{theorems,skins,hooks}
\newtcbtheorem[number within=chapter]{corollary}{Corollary}
{%
	enhanced
	,breakable
	,colback = myp!10
	,frame hidden
	,boxrule = 0sp
	,borderline west = {2pt}{0pt}{myp!85!black}
	,sharp corners
	,detach title
	,before upper = \tcbtitle\par\smallskip
	,coltitle = myp!85!black
	,fonttitle = \bfseries\sffamily
	,description font = \mdseries
	,separator sign none
	,segmentation style={solid, myp!85!black}
}
{th}


%================================
% LENMA
%================================

\tcbuselibrary{theorems,skins,hooks}
\newtcbtheorem[number within=section]{Lenma}{Lenma}
{%
	enhanced,
	breakable,
	colback = mylenmabg,
	frame hidden,
	boxrule = 0sp,
	borderline west = {2pt}{0pt}{mylenmafr},
	sharp corners,
	detach title,
	before upper = \tcbtitle\par\smallskip,
	coltitle = mylenmafr,
	fonttitle = \bfseries\sffamily,
	description font = \mdseries,
	separator sign none,
	segmentation style={solid, mylenmafr},
}
{th}

\tcbuselibrary{theorems,skins,hooks}
\newtcbtheorem[number within=chapter]{lenma}{Lenma}
{%
	enhanced,
	breakable,
	colback = mylenmabg,
	frame hidden,
	boxrule = 0sp,
	borderline west = {2pt}{0pt}{mylenmafr},
	sharp corners,
	detach title,
	before upper = \tcbtitle\par\smallskip,
	coltitle = mylenmafr,
	fonttitle = \bfseries\sffamily,
	description font = \mdseries,
	separator sign none,
	segmentation style={solid, mylenmafr},
}
{th}


%================================
% PROPOSITION
%================================

\tcbuselibrary{theorems,skins,hooks}
\newtcbtheorem[number within=section]{Prop}{Proposition}
{%
	enhanced,
	breakable,
	colback = mypropbg,
	frame hidden,
	boxrule = 0sp,
	borderline west = {2pt}{0pt}{mypropfr},
	sharp corners,
	detach title,
	before upper = \tcbtitle\par\smallskip,
	coltitle = mypropfr,
	fonttitle = \bfseries\sffamily,
	description font = \mdseries,
	separator sign none,
	segmentation style={solid, mypropfr},
}
{th}

\tcbuselibrary{theorems,skins,hooks}
\newtcbtheorem[number within=chapter]{prop}{Proposition}
{%
	enhanced,
	breakable,
	colback = mypropbg,
	frame hidden,
	boxrule = 0sp,
	borderline west = {2pt}{0pt}{mypropfr},
	sharp corners,
	detach title,
	before upper = \tcbtitle\par\smallskip,
	coltitle = mypropfr,
	fonttitle = \bfseries\sffamily,
	description font = \mdseries,
	separator sign none,
	segmentation style={solid, mypropfr},
}
{th}


%================================
% CLAIM
%================================

\tcbuselibrary{theorems,skins,hooks}
\newtcbtheorem[number within=section]{claim}{Claim}
{%
	enhanced
	,breakable
	,colback = myg!10
	,frame hidden
	,boxrule = 0sp
	,borderline west = {2pt}{0pt}{myg}
	,sharp corners
	,detach title
	,before upper = \tcbtitle\par\smallskip
	,coltitle = myg!85!black
	,fonttitle = \bfseries\sffamily
	,description font = \mdseries
	,separator sign none
	,segmentation style={solid, myg!85!black}
}
{th}



%================================
% Exercise
%================================

\tcbuselibrary{theorems,skins,hooks}
\newtcbtheorem[number within=section]{Exercise}{Exercise}
{%
	enhanced,
	breakable,
	colback = myexercisebg,
	frame hidden,
	boxrule = 0sp,
	borderline west = {2pt}{0pt}{myexercisefg},
	sharp corners,
	detach title,
	before upper = \tcbtitle\par\smallskip,
	coltitle = myexercisefg,
	fonttitle = \bfseries\sffamily,
	description font = \mdseries,
	separator sign none,
	segmentation style={solid, myexercisefg},
}
{th}

\tcbuselibrary{theorems,skins,hooks}
\newtcbtheorem[number within=chapter]{exercise}{Exercise}
{%
	enhanced,
	breakable,
	colback = myexercisebg,
	frame hidden,
	boxrule = 0sp,
	borderline west = {2pt}{0pt}{myexercisefg},
	sharp corners,
	detach title,
	before upper = \tcbtitle\par\smallskip,
	coltitle = myexercisefg,
	fonttitle = \bfseries\sffamily,
	description font = \mdseries,
	separator sign none,
	segmentation style={solid, myexercisefg},
}
{th}

%================================
% EXAMPLE BOX
%================================

\newtcbtheorem[number within=section]{Example}{Example}
{%
	colback = myexamplebg
	,breakable
	,colframe = myexamplefr
	,coltitle = myexampleti
	,boxrule = 1pt
	,sharp corners
	,detach title
	,before upper=\tcbtitle\par\smallskip
	,fonttitle = \bfseries
	,description font = \mdseries
	,separator sign none
	,description delimiters parenthesis
}
{ex}

\newtcbtheorem[number within=chapter]{example}{Example}
{%
	colback = myexamplebg
	,breakable
	,colframe = myexamplefr
	,coltitle = myexampleti
	,boxrule = 1pt
	,sharp corners
	,detach title
	,before upper=\tcbtitle\par\smallskip
	,fonttitle = \bfseries
	,description font = \mdseries
	,separator sign none
	,description delimiters parenthesis
}
{ex}

%================================
% DEFINITION BOX
%================================

\newtcbtheorem[number within=section]{Definition}{Definition}{enhanced,
	before skip=2mm,after skip=2mm, colback=red!5,colframe=red!80!black,boxrule=0.5mm,
	attach boxed title to top left={xshift=1cm,yshift*=1mm-\tcboxedtitleheight}, varwidth boxed title*=-3cm,
	boxed title style={frame code={
					\path[fill=tcbcolback]
					([yshift=-1mm,xshift=-1mm]frame.north west)
					arc[start angle=0,end angle=180,radius=1mm]
					([yshift=-1mm,xshift=1mm]frame.north east)
					arc[start angle=180,end angle=0,radius=1mm];
					\path[left color=tcbcolback!60!black,right color=tcbcolback!60!black,
						middle color=tcbcolback!80!black]
					([xshift=-2mm]frame.north west) -- ([xshift=2mm]frame.north east)
					[rounded corners=1mm]-- ([xshift=1mm,yshift=-1mm]frame.north east)
					-- (frame.south east) -- (frame.south west)
					-- ([xshift=-1mm,yshift=-1mm]frame.north west)
					[sharp corners]-- cycle;
				},interior engine=empty,
		},
	fonttitle=\bfseries,
	title={#2},#1}{def}
\newtcbtheorem[number within=chapter]{definition}{Definition}{enhanced,
	before skip=2mm,after skip=2mm, colback=red!5,colframe=red!80!black,boxrule=0.5mm,
	attach boxed title to top left={xshift=1cm,yshift*=1mm-\tcboxedtitleheight}, varwidth boxed title*=-3cm,
	boxed title style={frame code={
					\path[fill=tcbcolback]
					([yshift=-1mm,xshift=-1mm]frame.north west)
					arc[start angle=0,end angle=180,radius=1mm]
					([yshift=-1mm,xshift=1mm]frame.north east)
					arc[start angle=180,end angle=0,radius=1mm];
					\path[left color=tcbcolback!60!black,right color=tcbcolback!60!black,
						middle color=tcbcolback!80!black]
					([xshift=-2mm]frame.north west) -- ([xshift=2mm]frame.north east)
					[rounded corners=1mm]-- ([xshift=1mm,yshift=-1mm]frame.north east)
					-- (frame.south east) -- (frame.south west)
					-- ([xshift=-1mm,yshift=-1mm]frame.north west)
					[sharp corners]-- cycle;
				},interior engine=empty,
		},
	fonttitle=\bfseries,
	title={#2},#1}{def}



%================================
% Solution BOX
%================================

\makeatletter
\newtcbtheorem{question}{Question}{enhanced,
	breakable,
	colback=white,
	colframe=myb!80!black,
	attach boxed title to top left={yshift*=-\tcboxedtitleheight},
	fonttitle=\bfseries,
	title={#2},
	boxed title size=title,
	boxed title style={%
			sharp corners,
			rounded corners=northwest,
			colback=tcbcolframe,
			boxrule=0pt,
		},
	underlay boxed title={%
			\path[fill=tcbcolframe] (title.south west)--(title.south east)
			to[out=0, in=180] ([xshift=5mm]title.east)--
			(title.center-|frame.east)
			[rounded corners=\kvtcb@arc] |-
			(frame.north) -| cycle;
		},
	#1
}{def}
\makeatother

%================================
% SOLUTION BOX
%================================

\makeatletter
\newtcolorbox{solution}{enhanced,
	breakable,
	colback=white,
	colframe=myg!80!black,
	attach boxed title to top left={yshift*=-\tcboxedtitleheight},
	title=Solution,
	boxed title size=title,
	boxed title style={%
			sharp corners,
			rounded corners=northwest,
			colback=tcbcolframe,
			boxrule=0pt,
		},
	underlay boxed title={%
			\path[fill=tcbcolframe] (title.south west)--(title.south east)
			to[out=0, in=180] ([xshift=5mm]title.east)--
			(title.center-|frame.east)
			[rounded corners=\kvtcb@arc] |-
			(frame.north) -| cycle;
		},
}
\makeatother

%================================
% Question BOX
%================================

\makeatletter
\newtcbtheorem{qstion}{Question}{enhanced,
	breakable,
	colback=white,
	colframe=mygr,
	attach boxed title to top left={yshift*=-\tcboxedtitleheight},
	fonttitle=\bfseries,
	title={#2},
	boxed title size=title,
	boxed title style={%
			sharp corners,
			rounded corners=northwest,
			colback=tcbcolframe,
			boxrule=0pt,
		},
	underlay boxed title={%
			\path[fill=tcbcolframe] (title.south west)--(title.south east)
			to[out=0, in=180] ([xshift=5mm]title.east)--
			(title.center-|frame.east)
			[rounded corners=\kvtcb@arc] |-
			(frame.north) -| cycle;
		},
	#1
}{def}
\makeatother

\newtcbtheorem[number within=chapter]{wconc}{Wrong Concept}{
	breakable,
	enhanced,
	colback=white,
	colframe=myr,
	arc=0pt,
	outer arc=0pt,
	fonttitle=\bfseries\sffamily\large,
	colbacktitle=myr,
	attach boxed title to top left={},
	boxed title style={
			enhanced,
			skin=enhancedfirst jigsaw,
			arc=3pt,
			bottom=0pt,
			interior style={fill=myr}
		},
	#1
}{def}



%================================
% NOTE BOX
%================================

\usetikzlibrary{arrows,calc,shadows.blur}
\tcbuselibrary{skins}
\newtcolorbox{note}[1][]{%
	enhanced jigsaw,
	colback=gray!20!white,%
	colframe=gray!80!black,
	size=small,
	boxrule=1pt,
	title=\textbf{Note:},
	halign title=flush center,
	coltitle=black,
	breakable,
	drop shadow=black!50!white,
	attach boxed title to top left={xshift=1cm,yshift=-\tcboxedtitleheight/2,yshifttext=-\tcboxedtitleheight/2},
	minipage boxed title=1.5cm,
	boxed title style={%
			colback=white,
			size=fbox,
			boxrule=1pt,
			boxsep=2pt,
			underlay={%
					\coordinate (dotA) at ($(interior.west) + (-0.5pt,0)$);
					\coordinate (dotB) at ($(interior.east) + (0.5pt,0)$);
					\begin{scope}
						\clip (interior.north west) rectangle ([xshift=3ex]interior.east);
						\filldraw [white, blur shadow={shadow opacity=60, shadow yshift=-.75ex}, rounded corners=2pt] (interior.north west) rectangle (interior.south east);
					\end{scope}
					\begin{scope}[gray!80!black]
						\fill (dotA) circle (2pt);
						\fill (dotB) circle (2pt);
					\end{scope}
				},
		},
	#1,
}

%%%%%%%%%%%%%%%%%%%%%%%%%%%%%%
% SELF MADE COMMANDS
%%%%%%%%%%%%%%%%%%%%%%%%%%%%%%


\newcommand{\thm}[2]{\begin{Theorem}{#1}{}#2\end{Theorem}}
\newcommand{\cor}[2]{\begin{Corollary}{#1}{}#2\end{Corollary}}
\newcommand{\mlenma}[2]{\begin{Lenma}{#1}{}#2\end{Lenma}}
\newcommand{\mprop}[2]{\begin{Prop}{#1}{}#2\end{Prop}}
\newcommand{\clm}[3]{\begin{claim}{#1}{#2}#3\end{claim}}
\newcommand{\wc}[2]{\begin{wconc}{#1}{}\setlength{\parindent}{1cm}#2\end{wconc}}
\newcommand{\thmcon}[1]{\begin{Theoremcon}{#1}\end{Theoremcon}}
\newcommand{\ex}[2]{\begin{Example}{#1}{}#2\end{Example}}
\newcommand{\dfn}[2]{\begin{Definition}[colbacktitle=red!75!black]{#1}{}#2\end{Definition}}
\newcommand{\dfnc}[2]{\begin{definition}[colbacktitle=red!75!black]{#1}{}#2\end{definition}}
\newcommand{\qs}[2]{\begin{question}{#1}{}#2\end{question}}
\newcommand{\pf}[2]{\begin{myproof}[#1]#2\end{myproof}}
\newcommand{\nt}[1]{\begin{note}#1\end{note}}

\newcommand*\circled[1]{\tikz[baseline=(char.base)]{
		\node[shape=circle,draw,inner sep=1pt] (char) {#1};}}
\newcommand\getcurrentref[1]{%
	\ifnumequal{\value{#1}}{0}
	{??}
	{\the\value{#1}}%
}
\newcommand{\getCurrentSectionNumber}{\getcurrentref{section}}
\newenvironment{myproof}[1][\proofname]{%
	\proof[\bfseries #1: ]%
}{\endproof}

\newcommand{\mclm}[2]{\begin{myclaim}[#1]#2\end{myclaim}}
\newenvironment{myclaim}[1][\claimname]{\proof[\bfseries #1: ]}{}

\newcounter{mylabelcounter}

\makeatletter
\newcommand{\setword}[2]{%
	\phantomsection
	#1\def\@currentlabel{\unexpanded{#1}}\label{#2}%
}
\makeatother




\tikzset{
	symbol/.style={
			draw=none,
			every to/.append style={
					edge node={node [sloped, allow upside down, auto=false]{$#1$}}}
		}
}


% deliminators
\DeclarePairedDelimiter{\abs}{\lvert}{\rvert}
\DeclarePairedDelimiter{\norm}{\lVert}{\rVert}

\DeclarePairedDelimiter{\ceil}{\lceil}{\rceil}
\DeclarePairedDelimiter{\floor}{\lfloor}{\rfloor}
\DeclarePairedDelimiter{\round}{\lfloor}{\rceil}

\newsavebox\diffdbox
\newcommand{\slantedromand}{{\mathpalette\makesl{d}}}
\newcommand{\makesl}[2]{%
\begingroup
\sbox{\diffdbox}{$\mathsurround=0pt#1\mathrm{#2}$}%
\pdfsave
\pdfsetmatrix{1 0 0.2 1}%
\rlap{\usebox{\diffdbox}}%
\pdfrestore
\hskip\wd\diffdbox
\endgroup
}
\newcommand{\dd}[1][]{\ensuremath{\mathop{}\!\ifstrempty{#1}{%
\slantedromand\@ifnextchar^{\hspace{0.2ex}}{\hspace{0.1ex}}}%
{\slantedromand\hspace{0.2ex}^{#1}}}}
\ProvideDocumentCommand\dv{o m g}{%
  \ensuremath{%
    \IfValueTF{#3}{%
      \IfNoValueTF{#1}{%
        \frac{\dd #2}{\dd #3}%
      }{%
        \frac{\dd^{#1} #2}{\dd #3^{#1}}%
      }%
    }{%
      \IfNoValueTF{#1}{%
        \frac{\dd}{\dd #2}%
      }{%
        \frac{\dd^{#1}}{\dd #2^{#1}}%
      }%
    }%
  }%
}
\providecommand*{\pdv}[3][]{\frac{\partial^{#1}#2}{\partial#3^{#1}}}
%  - others
\DeclareMathOperator{\Lap}{\mathcal{L}}
\DeclareMathOperator{\Var}{Var} % varience
\DeclareMathOperator{\Cov}{Cov} % covarience
\DeclareMathOperator{\E}{E} % expected

% Since the amsthm package isn't loaded

% I prefer the slanted \leq
\let\oldleq\leq % save them in case they're every wanted
\let\oldgeq\geq
\renewcommand{\leq}{\leqslant}
\renewcommand{\geq}{\geqslant}

% % redefine matrix env to allow for alignment, use r as default
% \renewcommand*\env@matrix[1][r]{\hskip -\arraycolsep
%     \let\@ifnextchar\new@ifnextchar
%     \array{*\c@MaxMatrixCols #1}}


%\usepackage{framed}
%\usepackage{titletoc}
%\usepackage{etoolbox}
%\usepackage{lmodern}


%\patchcmd{\tableofcontents}{\contentsname}{\sffamily\contentsname}{}{}

%\renewenvironment{leftbar}
%{\def\FrameCommand{\hspace{6em}%
%		{\color{myyellow}\vrule width 2pt depth 6pt}\hspace{1em}}%
%	\MakeFramed{\parshape 1 0cm \dimexpr\textwidth-6em\relax\FrameRestore}\vskip2pt%
%}
%{\endMakeFramed}

%\titlecontents{chapter}
%[0em]{\vspace*{2\baselineskip}}
%{\parbox{4.5em}{%
%		\hfill\Huge\sffamily\bfseries\color{myred}\thecontentspage}%
%	\vspace*{-2.3\baselineskip}\leftbar\textsc{\small\chaptername~\thecontentslabel}\\\sffamily}
%{}{\endleftbar}
%\titlecontents{section}
%[8.4em]
%{\sffamily\contentslabel{3em}}{}{}
%{\hspace{0.5em}\nobreak\itshape\color{myred}\contentspage}
%\titlecontents{subsection}
%[8.4em]
%{\sffamily\contentslabel{3em}}{}{}  
%{\hspace{0.5em}\nobreak\itshape\color{myred}\contentspage}



%%%%%%%%%%%%%%%%%%%%%%%%%%%%%%%%%%%%%%%%%%%
% TABLE OF CONTENTS
%%%%%%%%%%%%%%%%%%%%%%%%%%%%%%%%%%%%%%%%%%%

\usepackage{tikz}
\definecolor{doc}{RGB}{0,60,110}
\usepackage{titletoc}
\contentsmargin{0cm}
\titlecontents{chapter}[3.7pc]
{\addvspace{30pt}%
	\begin{tikzpicture}[remember picture, overlay]%
		\draw[fill=doc!60,draw=doc!60] (-7,-.1) rectangle (-0.9,.5);%
		\pgftext[left,x=-3.5cm,y=0.2cm]{\color{white}\Large\sc\bfseries Chapter\ \thecontentslabel};%
	\end{tikzpicture}\color{doc!60}\large\sc\bfseries}%
{}
{}
{\;\titlerule\;\large\sc\bfseries Page \thecontentspage
	\begin{tikzpicture}[remember picture, overlay]
		\draw[fill=doc!60,draw=doc!60] (2pt,0) rectangle (4,0.1pt);
	\end{tikzpicture}}%
\titlecontents{section}[3.7pc]
{\addvspace{2pt}}
{\contentslabel[\thecontentslabel]{2pc}}
{}
{\hfill\small \thecontentspage}
[]
\titlecontents*{subsection}[3.7pc]
{\addvspace{-1pt}\small}
{}
{}
{\ --- \small\thecontentspage}
[ \textbullet\ ][]

\makeatletter
\renewcommand{\tableofcontents}{%
	\chapter*{%
	  \vspace*{-20\p@}%
	  \begin{tikzpicture}[remember picture, overlay]%
		  \pgftext[right,x=15cm,y=0.2cm]{\color{doc!60}\Huge\sc\bfseries \contentsname};%
		  \draw[fill=doc!60,draw=doc!60] (13,-.75) rectangle (20,1);%
		  \clip (13,-.75) rectangle (20,1);
		  \pgftext[right,x=15cm,y=0.2cm]{\color{white}\Huge\sc\bfseries \contentsname};%
	  \end{tikzpicture}}%
	\@starttoc{toc}}
\makeatother


% Comandi di Bastiality %
\newcommand{\innerproduct}[2]{\langle #1, #2 \rangle}
\newcommand{\generators}[1]{\langle #1 \rangle}
\newcommand{\powerset}[0]{\mathcal{P}}
\newcommand{\indip}{\perp \!\!\! \perp}

% Comandi di GioLaPalma %
\newcommand{\der}[0]{\implies^*}
\newcommand{\red}[1]{\textcolor{red}{#1}}
\newcommand{\esempio}[1]{\ex{}{#1}}
\newcommand{\teorema}[1]{\thm{}{#1}}
\newcommand{\dimostrazione}[1]{\pf{Dimostrazione}{#1}}
\setlength{\parindent}{0pt}
\newcommand{\osservazione}[1]{\nt{#1}}
\newcommand{\inside}[1]{
  \langle #1 \rangle
}
      

%From M275 "Topology" at SJSU
\newcommand{\id}{\mathrm{id}}
\newcommand{\taking}[1]{\xrightarrow{#1}}
\newcommand{\inv}{^{-1}}

%From M170 "Introduction to Graph Theory" at SJSU
\DeclareMathOperator{\diam}{diam}
\DeclareMathOperator{\ord}{ord}
\newcommand{\defeq}{\overset{\mathrm{def}}{=}}

%From the USAMO .tex files
\newcommand{\ts}{\textsuperscript}
\newcommand{\dg}{^\circ}
\newcommand{\ii}{\item}

% % From Math 55 and Math 145 at Harvard
% \newenvironment{subproof}[1][Proof]{%
% \begin{proof}[#1] \renewcommand{\qedsymbol}{$\blacksquare$}}%
% {\end{proof}}

\newcommand{\liff}{\leftrightarrow}
\newcommand{\lthen}{\rightarrow}
\newcommand{\opname}{\operatorname}
\newcommand{\surjto}{\twoheadrightarrow}
\newcommand{\injto}{\hookrightarrow}
\newcommand{\On}{\mathrm{On}} % ordinals
\DeclareMathOperator{\img}{im} % Image
\DeclareMathOperator{\Img}{Im} % Image
\DeclareMathOperator{\coker}{coker} % Cokernel
\DeclareMathOperator{\Coker}{Coker} % Cokernel
\DeclareMathOperator{\Ker}{Ker} % Kernel
\DeclareMathOperator{\rank}{rank}
\DeclareMathOperator{\Spec}{Spec} % spectrum
\DeclareMathOperator{\Tr}{Tr} % trace
\DeclareMathOperator{\pr}{pr} % projection
\DeclareMathOperator{\ext}{ext} % extension
\DeclareMathOperator{\pred}{pred} % predecessor
\DeclareMathOperator{\dom}{dom} % domain
\DeclareMathOperator{\ran}{ran} % range
\DeclareMathOperator{\Hom}{Hom} % homomorphism
\DeclareMathOperator{\Mor}{Mor} % morphisms
\DeclareMathOperator{\End}{End} % endomorphism

\newcommand{\eps}{\epsilon}
\newcommand{\veps}{\varepsilon}
\newcommand{\ol}{\overline}
\newcommand{\ul}{\underline}
\newcommand{\wt}{\widetilde}
\newcommand{\wh}{\widehat}
\newcommand{\vocab}[1]{\textbf{\color{blue} #1}}
\providecommand{\half}{\frac{1}{2}}
\newcommand{\dang}{\measuredangle} %% Directed angle
\newcommand{\ray}[1]{\overrightarrow{#1}}
\newcommand{\seg}[1]{\overline{#1}}
\newcommand{\arc}[1]{\wideparen{#1}}
\DeclareMathOperator{\cis}{cis}
\DeclareMathOperator*{\lcm}{lcm}
\DeclareMathOperator*{\argmin}{arg min}
\DeclareMathOperator*{\argmax}{arg max}
\newcommand{\cycsum}{\sum_{\mathrm{cyc}}}
\newcommand{\symsum}{\sum_{\mathrm{sym}}}
\newcommand{\cycprod}{\prod_{\mathrm{cyc}}}
\newcommand{\symprod}{\prod_{\mathrm{sym}}}
\newcommand{\Qed}{\begin{flushright}\qed\end{flushright}}
\newcommand{\parinn}{\setlength{\parindent}{1cm}}
\newcommand{\parinf}{\setlength{\parindent}{0cm}}
% \newcommand{\norm}{\|\cdot\|}
\newcommand{\inorm}{\norm_{\infty}}
\newcommand{\opensets}{\{V_{\alpha}\}_{\alpha\in I}}
\newcommand{\oset}{V_{\alpha}}
\newcommand{\opset}[1]{V_{\alpha_{#1}}}
\newcommand{\lub}{\text{lub}}
\newcommand{\del}[2]{\frac{\partial #1}{\partial #2}}
\newcommand{\Del}[3]{\frac{\partial^{#1} #2}{\partial^{#1} #3}}
\newcommand{\deld}[2]{\dfrac{\partial #1}{\partial #2}}
\newcommand{\Deld}[3]{\dfrac{\partial^{#1} #2}{\partial^{#1} #3}}
\newcommand{\lm}{\lambda}
\newcommand{\uin}{\mathbin{\rotatebox[origin=c]{90}{$\in$}}}
\newcommand{\usubset}{\mathbin{\rotatebox[origin=c]{90}{$\subset$}}}
\newcommand{\lt}{\left}
\newcommand{\rt}{\right}
\newcommand{\bs}[1]{\boldsymbol{#1}}
\newcommand{\exs}{\exists}
\newcommand{\st}{\strut}
\newcommand{\dps}[1]{\displaystyle{#1}}

\newcommand{\sol}{\setlength{\parindent}{0cm}\textbf{\textit{Solution:}}\setlength{\parindent}{1cm} }
\newcommand{\solve}[1]{\setlength{\parindent}{0cm}\textbf{\textit{Solution: }}\setlength{\parindent}{1cm}#1 \Qed}

% Things Lie
\newcommand{\kb}{\mathfrak b}
\newcommand{\kg}{\mathfrak g}
\newcommand{\kh}{\mathfrak h}
\newcommand{\kn}{\mathfrak n}
\newcommand{\ku}{\mathfrak u}
\newcommand{\kz}{\mathfrak z}
\DeclareMathOperator{\Ext}{Ext} % Ext functor
\DeclareMathOperator{\Tor}{Tor} % Tor functor
\newcommand{\gl}{\opname{\mathfrak{gl}}} % frak gl group
\renewcommand{\sl}{\opname{\mathfrak{sl}}} % frak sl group chktex 6

% More script letters etc.
\newcommand{\SA}{\mathcal A}
\newcommand{\SB}{\mathcal B}
\newcommand{\SC}{\mathcal C}
\newcommand{\SF}{\mathcal F}
\newcommand{\SG}{\mathcal G}
\newcommand{\SH}{\mathcal H}
\newcommand{\OO}{\mathcal O}

\newcommand{\SCA}{\mathscr A}
\newcommand{\SCB}{\mathscr B}
\newcommand{\SCC}{\mathscr C}
\newcommand{\SCD}{\mathscr D}
\newcommand{\SCE}{\mathscr E}
\newcommand{\SCF}{\mathscr F}
\newcommand{\SCG}{\mathscr G}
\newcommand{\SCH}{\mathscr H}

% Mathfrak primes
\newcommand{\km}{\mathfrak m}
\newcommand{\kp}{\mathfrak p}
\newcommand{\kq}{\mathfrak q}

% number sets
\newcommand{\RR}[1][]{\ensuremath{\ifstrempty{#1}{\mathbb{R}}{\mathbb{R}^{#1}}}}
\newcommand{\NN}[1][]{\ensuremath{\ifstrempty{#1}{\mathbb{N}}{\mathbb{N}^{#1}}}}
\newcommand{\ZZ}[1][]{\ensuremath{\ifstrempty{#1}{\mathbb{Z}}{\mathbb{Z}^{#1}}}}
\newcommand{\QQ}[1][]{\ensuremath{\ifstrempty{#1}{\mathbb{Q}}{\mathbb{Q}^{#1}}}}
\newcommand{\CC}[1][]{\ensuremath{\ifstrempty{#1}{\mathbb{C}}{\mathbb{C}^{#1}}}}
\newcommand{\PP}[1][]{\ensuremath{\ifstrempty{#1}{\mathbb{P}}{\mathbb{P}^{#1}}}}
\newcommand{\HH}[1][]{\ensuremath{\ifstrempty{#1}{\mathbb{H}}{\mathbb{H}^{#1}}}}
\newcommand{\FF}[1][]{\ensuremath{\ifstrempty{#1}{\mathbb{F}}{\mathbb{F}^{#1}}}}
% expected value
\newcommand{\EE}{\ensuremath{\mathbb{E}}}
\newcommand{\charin}{\text{ char }}
\DeclareMathOperator{\sign}{sign}
\DeclareMathOperator{\Aut}{Aut}
\DeclareMathOperator{\Inn}{Inn}
\DeclareMathOperator{\Syl}{Syl}
\DeclareMathOperator{\Gal}{Gal}
\DeclareMathOperator{\GL}{GL} % General linear group
\DeclareMathOperator{\SL}{SL} % Special linear group

%---------------------------------------
% BlackBoard Math Fonts :-
%---------------------------------------

%Captital Letters
\newcommand{\bbA}{\mathbb{A}}	\newcommand{\bbB}{\mathbb{B}}
\newcommand{\bbC}{\mathbb{C}}	\newcommand{\bbD}{\mathbb{D}}
\newcommand{\bbE}{\mathbb{E}}	\newcommand{\bbF}{\mathbb{F}}
\newcommand{\bbG}{\mathbb{G}}	\newcommand{\bbH}{\mathbb{H}}
\newcommand{\bbI}{\mathbb{I}}	\newcommand{\bbJ}{\mathbb{J}}
\newcommand{\bbK}{\mathbb{K}}	\newcommand{\bbL}{\mathbb{L}}
\newcommand{\bbM}{\mathbb{M}}	\newcommand{\bbN}{\mathbb{N}}
\newcommand{\bbO}{\mathbb{O}}	\newcommand{\bbP}{\mathbb{P}}
\newcommand{\bbQ}{\mathbb{Q}}	\newcommand{\bbR}{\mathbb{R}}
\newcommand{\bbS}{\mathbb{S}}	\newcommand{\bbT}{\mathbb{T}}
\newcommand{\bbU}{\mathbb{U}}	\newcommand{\bbV}{\mathbb{V}}
\newcommand{\bbW}{\mathbb{W}}	\newcommand{\bbX}{\mathbb{X}}
\newcommand{\bbY}{\mathbb{Y}}	\newcommand{\bbZ}{\mathbb{Z}}

%---------------------------------------
% MathCal Fonts :-
%---------------------------------------

%Captital Letters
\newcommand{\mcA}{\mathcal{A}}	\newcommand{\mcB}{\mathcal{B}}
\newcommand{\mcC}{\mathcal{C}}	\newcommand{\mcD}{\mathcal{D}}
\newcommand{\mcE}{\mathcal{E}}	\newcommand{\mcF}{\mathcal{F}}
\newcommand{\mcG}{\mathcal{G}}	\newcommand{\mcH}{\mathcal{H}}
\newcommand{\mcI}{\mathcal{I}}	\newcommand{\mcJ}{\mathcal{J}}
\newcommand{\mcK}{\mathcal{K}}	\newcommand{\mcL}{\mathcal{L}}
\newcommand{\mcM}{\mathcal{M}}	\newcommand{\mcN}{\mathcal{N}}
\newcommand{\mcO}{\mathcal{O}}	\newcommand{\mcP}{\mathcal{P}}
\newcommand{\mcQ}{\mathcal{Q}}	\newcommand{\mcR}{\mathcal{R}}
\newcommand{\mcS}{\mathcal{S}}	\newcommand{\mcT}{\mathcal{T}}
\newcommand{\mcU}{\mathcal{U}}	\newcommand{\mcV}{\mathcal{V}}
\newcommand{\mcW}{\mathcal{W}}	\newcommand{\mcX}{\mathcal{X}}
\newcommand{\mcY}{\mathcal{Y}}	\newcommand{\mcZ}{\mathcal{Z}}


%---------------------------------------
% Bold Math Fonts :-
%---------------------------------------

%Captital Letters
\newcommand{\bmA}{\boldsymbol{A}}	\newcommand{\bmB}{\boldsymbol{B}}
\newcommand{\bmC}{\boldsymbol{C}}	\newcommand{\bmD}{\boldsymbol{D}}
\newcommand{\bmE}{\boldsymbol{E}}	\newcommand{\bmF}{\boldsymbol{F}}
\newcommand{\bmG}{\boldsymbol{G}}	\newcommand{\bmH}{\boldsymbol{H}}
\newcommand{\bmI}{\boldsymbol{I}}	\newcommand{\bmJ}{\boldsymbol{J}}
\newcommand{\bmK}{\boldsymbol{K}}	\newcommand{\bmL}{\boldsymbol{L}}
\newcommand{\bmM}{\boldsymbol{M}}	\newcommand{\bmN}{\boldsymbol{N}}
\newcommand{\bmO}{\boldsymbol{O}}	\newcommand{\bmP}{\boldsymbol{P}}
\newcommand{\bmQ}{\boldsymbol{Q}}	\newcommand{\bmR}{\boldsymbol{R}}
\newcommand{\bmS}{\boldsymbol{S}}	\newcommand{\bmT}{\boldsymbol{T}}
\newcommand{\bmU}{\boldsymbol{U}}	\newcommand{\bmV}{\boldsymbol{V}}
\newcommand{\bmW}{\boldsymbol{W}}	\newcommand{\bmX}{\boldsymbol{X}}
\newcommand{\bmY}{\boldsymbol{Y}}	\newcommand{\bmZ}{\boldsymbol{Z}}
%Small Letters
\newcommand{\bma}{\boldsymbol{a}}	\newcommand{\bmb}{\boldsymbol{b}}
\newcommand{\bmc}{\boldsymbol{c}}	\newcommand{\bmd}{\boldsymbol{d}}
\newcommand{\bme}{\boldsymbol{e}}	\newcommand{\bmf}{\boldsymbol{f}}
\newcommand{\bmg}{\boldsymbol{g}}	\newcommand{\bmh}{\boldsymbol{h}}
\newcommand{\bmi}{\boldsymbol{i}}	\newcommand{\bmj}{\boldsymbol{j}}
\newcommand{\bmk}{\boldsymbol{k}}	\newcommand{\bml}{\boldsymbol{l}}
\newcommand{\bmm}{\boldsymbol{m}}	\newcommand{\bmn}{\boldsymbol{n}}
\newcommand{\bmo}{\boldsymbol{o}}	\newcommand{\bmp}{\boldsymbol{p}}
\newcommand{\bmq}{\boldsymbol{q}}	\newcommand{\bmr}{\boldsymbol{r}}
\newcommand{\bms}{\boldsymbol{s}}	\newcommand{\bmt}{\boldsymbol{t}}
\newcommand{\bmu}{\boldsymbol{u}}	\newcommand{\bmv}{\boldsymbol{v}}
\newcommand{\bmw}{\boldsymbol{w}}	\newcommand{\bmx}{\boldsymbol{x}}
\newcommand{\bmy}{\boldsymbol{y}}	\newcommand{\bmz}{\boldsymbol{z}}

%---------------------------------------
% Scr Math Fonts :-
%---------------------------------------

\newcommand{\sA}{{\mathscr{A}}}   \newcommand{\sB}{{\mathscr{B}}}
\newcommand{\sC}{{\mathscr{C}}}   \newcommand{\sD}{{\mathscr{D}}}
\newcommand{\sE}{{\mathscr{E}}}   \newcommand{\sF}{{\mathscr{F}}}
\newcommand{\sG}{{\mathscr{G}}}   \newcommand{\sH}{{\mathscr{H}}}
\newcommand{\sI}{{\mathscr{I}}}   \newcommand{\sJ}{{\mathscr{J}}}
\newcommand{\sK}{{\mathscr{K}}}   \newcommand{\sL}{{\mathscr{L}}}
\newcommand{\sM}{{\mathscr{M}}}   \newcommand{\sN}{{\mathscr{N}}}
\newcommand{\sO}{{\mathscr{O}}}   \newcommand{\sP}{{\mathscr{P}}}
\newcommand{\sQ}{{\mathscr{Q}}}   \newcommand{\sR}{{\mathscr{R}}}
\newcommand{\sS}{{\mathscr{S}}}   \newcommand{\sT}{{\mathscr{T}}}
\newcommand{\sU}{{\mathscr{U}}}   \newcommand{\sV}{{\mathscr{V}}}
\newcommand{\sW}{{\mathscr{W}}}   \newcommand{\sX}{{\mathscr{X}}}
\newcommand{\sY}{{\mathscr{Y}}}   \newcommand{\sZ}{{\mathscr{Z}}}


%---------------------------------------
% Math Fraktur Font
%---------------------------------------

%Captital Letters
\newcommand{\mfA}{\mathfrak{A}}	\newcommand{\mfB}{\mathfrak{B}}
\newcommand{\mfC}{\mathfrak{C}}	\newcommand{\mfD}{\mathfrak{D}}
\newcommand{\mfE}{\mathfrak{E}}	\newcommand{\mfF}{\mathfrak{F}}
\newcommand{\mfG}{\mathfrak{G}}	\newcommand{\mfH}{\mathfrak{H}}
\newcommand{\mfI}{\mathfrak{I}}	\newcommand{\mfJ}{\mathfrak{J}}
\newcommand{\mfK}{\mathfrak{K}}	\newcommand{\mfL}{\mathfrak{L}}
\newcommand{\mfM}{\mathfrak{M}}	\newcommand{\mfN}{\mathfrak{N}}
\newcommand{\mfO}{\mathfrak{O}}	\newcommand{\mfP}{\mathfrak{P}}
\newcommand{\mfQ}{\mathfrak{Q}}	\newcommand{\mfR}{\mathfrak{R}}
\newcommand{\mfS}{\mathfrak{S}}	\newcommand{\mfT}{\mathfrak{T}}
\newcommand{\mfU}{\mathfrak{U}}	\newcommand{\mfV}{\mathfrak{V}}
\newcommand{\mfW}{\mathfrak{W}}	\newcommand{\mfX}{\mathfrak{X}}
\newcommand{\mfY}{\mathfrak{Y}}	\newcommand{\mfZ}{\mathfrak{Z}}
%Small Letters
\newcommand{\mfa}{\mathfrak{a}}	\newcommand{\mfb}{\mathfrak{b}}
\newcommand{\mfc}{\mathfrak{c}}	\newcommand{\mfd}{\mathfrak{d}}
\newcommand{\mfe}{\mathfrak{e}}	\newcommand{\mff}{\mathfrak{f}}
\newcommand{\mfg}{\mathfrak{g}}	\newcommand{\mfh}{\mathfrak{h}}
\newcommand{\mfi}{\mathfrak{i}}	\newcommand{\mfj}{\mathfrak{j}}
\newcommand{\mfk}{\mathfrak{k}}	\newcommand{\mfl}{\mathfrak{l}}
\newcommand{\mfm}{\mathfrak{m}}	\newcommand{\mfn}{\mathfrak{n}}
\newcommand{\mfo}{\mathfrak{o}}	\newcommand{\mfp}{\mathfrak{p}}
\newcommand{\mfq}{\mathfrak{q}}	\newcommand{\mfr}{\mathfrak{r}}
\newcommand{\mfs}{\mathfrak{s}}	\newcommand{\mft}{\mathfrak{t}}
\newcommand{\mfu}{\mathfrak{u}}	\newcommand{\mfv}{\mathfrak{v}}
\newcommand{\mfw}{\mathfrak{w}}	\newcommand{\mfx}{\mathfrak{x}}
\newcommand{\mfy}{\mathfrak{y}}	\newcommand{\mfz}{\mathfrak{z}}


\usepackage[T1]{fontenc}
\usepackage[utf8]{inputenc}
\usepackage{xcolor}
\usepackage{listings}
\usepackage{array}
\usepackage{longtable}

%==============================================================
% Stile per i blocchi di codice (C, C++, Java, Python)
%==============================================================
\definecolor{codebg}{HTML}{F7F7FB}
\definecolor{codekw}{HTML}{2D6FB1}
\definecolor{codecm}{HTML}{388C46}
\definecolor{codestr}{HTML}{C74440}
\definecolor{coderule}{HTML}{CCCCCC}

\lstset{
    basicstyle=\ttfamily\footnotesize,
    keywordstyle=\color{codekw}\bfseries,
    commentstyle=\color{codecm}\itshape,
    stringstyle=\color{codestr},
    numberstyle=\tiny\color{gray},
    showstringspaces=false,
    breaklines=true,
    breakatwhitespace=false,
    frame=single,
    rulecolor=\color{coderule},
    backgroundcolor=\color{codebg},
    tabsize=2,
    columns=fullflexible,
    keepspaces=true,
    captionpos=b,
    xleftmargin=1.2em,
    aboveskip=0.8em,
    belowskip=0.8em,
}

% codice inline
\newcommand{\code}[1]{\texttt{\small #1}}
% metodi dunder di Python: \dunder{init} -> __init__
\newcommand{\dunder}[1]{\texttt{\small \_\_#1\_\_}}

%==============================================================
% Box DOMANDA e SOLUZIONE
%   \begin{domanda}{A1}{titolo} ... \end{domanda}
%   \begin{soluzione}{A1 --- titolo} ... \end{soluzione}
% (sono ambienti, non macro: dentro ci si puo' mettere lstlisting)
%==============================================================
\newtcolorbox{domanda}[2]{%
    enhanced,
    breakable,
    colback=white,
    colframe=myb!80!black,
    fonttitle=\bfseries\sffamily,
    coltitle=white,
    colbacktitle=myb!80!black,
    title={#1 \quad\textbullet\quad #2},
    arc=2pt,
    left=3mm, right=3mm,
    boxrule=1pt,
}

\newtcolorbox{soluzione}[1]{%
    enhanced,
    breakable,
    colback=white,
    colframe=myg!80!black,
    fonttitle=\bfseries\sffamily,
    coltitle=white,
    colbacktitle=myg!80!black,
    title={Soluzione #1},
    arc=2pt,
    left=3mm, right=3mm,
    boxrule=1pt,
}

% riga finale della domanda con il rimando alla soluzione
\newcommand{\vaisol}[1]{%
    \par\vspace{4pt}\hfill{\footnotesize\itshape\color{gray}%
    Soluzione a pag.~\pageref{sol:#1}}%
}

\graphicspath{{./img/}}

%==============================================================
\begin{document}

\begin{titlepage}
    \centering
    \vspace*{2.5cm}
    {\Huge\sffamily\bfseries Overriding e funzioni virtuali\par}
    \vspace{0.5cm}
    {\LARGE\sffamily Domande d'esame avanzate\par}
    \vspace{1.0cm}
    {\LARGE\sffamily C \quad$\cdot$\quad C++ \quad$\cdot$\quad Java \quad$\cdot$\quad Python\par}
    \vspace{1.4cm}
    {\Large Appunti di Ingegneria del Software\par}
    \vspace{0.4cm}
    {\large 31 domande con soluzione ragionata\par}
    \vspace{2cm}
    {\large\itshape ``All'esame chiede piu' di quello che c'e' scritto:\\
    soprattutto le \textbf{differenze} tra i linguaggi.''\par}
    \vfill
    {\large Anno 3\par}
\end{titlepage}

\tableofcontents
\newpage

%==============================================================
\chapter{Come usare questo fascicolo}
%==============================================================

Questo fascicolo raccoglie \textbf{31 domande di difficolta' alta} su
\emph{overriding} e \emph{dispatch dinamico}, pensate per il tipo di
interrogazione in cui il professore parte da uno snippet di codice e chiede
\emph{``cosa stampa e perche'''}, oppure chiede di confrontare la stessa
feature nei quattro linguaggi del corso.

La struttura e' la seguente:

\begin{itemize}
    \item i capitoli \ref{ch:cpp}--\ref{ch:trasv} contengono \emph{solo le
          domande}, raggruppate per linguaggio (C++, Java, Python, C) piu' un
          gruppo trasversale: si possono usare come simulazione a libro chiuso;
    \item il capitolo \ref{ch:sol} contiene le \emph{soluzioni ragionate},
          nello stesso ordine; ogni domanda riporta in basso a destra il numero
          di pagina della propria soluzione.
\end{itemize}

I codici delle domande (\code{A1}, \code{B3}, \code{C5}, \dots) sono stabili:
la lettera indica il gruppo, il numero la posizione.

\nt{
    Tre idee fanno da filo conduttore a quasi tutte le domande, e conviene
    tenerle a mente perche' sono quelle su cui si sbaglia:
    \textbf{(1)} \emph{cosa} viene deciso a tempo di compilazione (overload,
    campi, argomenti di default, controllo di accesso) e cosa a tempo di
    esecuzione (solo il \emph{corpo} del metodo virtuale);
    \textbf{(2)} \emph{quando} il tipo dinamico di un oggetto e' gia' quello
    finale (in Java sempre, in C++ no durante costruzione e distruzione);
    \textbf{(3)} \emph{dove} avviene la ricerca del metodo (vtable a indice
    fisso in C++/Java, dizionario lungo l'MRO in Python, struct di puntatori a
    funzione scritta a mano in C).
}

\section{Fonti segnalate dal docente}

Il docente ha indicato come riferimento due voci di Wikipedia:

\begin{itemize}
    \item \textbf{Virtual inheritance} ---
          \texttt{\small en.wikipedia.org/wiki/Virtual\_inheritance}\\
          il problema del diamante con l'esempio \emph{Animal / Mammal /
          WingedAnimal / Bat}, la regola di inizializzazione della base virtuale
          e il costo in memoria. Corrisponde alle domande A4 e A9--A12.
    \item \textbf{Virtual function} ---
          \texttt{\small en.wikipedia.org/wiki/Virtual\_function}\\ dispatch
          dinamico e late binding, metodi virtuali puri e classi astratte,
          distruttori virtuali, chiamate virtuali durante la costruzione, vtable
          e relativo overhead, quali linguaggi hanno i metodi virtuali per
          default. Corrisponde alle domande A1, A8, A11, D1, E1 ed E2.
\end{itemize}

Le due voci sono il minimo sindacale: quasi tutte le domande di questo fascicolo
partono da li' e chiedono il passo successivo, cioe' \emph{perche'} il
meccanismo e' fatto cosi' e \emph{come cambia} negli altri tre linguaggi.

%==============================================================
\chapter{C++}\label{ch:cpp}
%==============================================================

\begin{domanda}{A1}{Chiamate virtuali in costruttore e distruttore}
Cosa stampa il programma? Motiva riga per riga.
\begin{lstlisting}[language=C++]
struct A {
    A() { std::cout << "A() "; f(); }
    virtual ~A() { g(); }
    virtual void f() { std::cout << "A::f "; }
    virtual void g() { std::cout << "A::g "; }
};
struct B : A {
    int x = 42;
    void f() override { std::cout << "B::f " << x << " "; }
    void g() override { std::cout << "B::g "; }
};
int main() { B b; }
\end{lstlisting}
E se in \code{A} avessi scritto \code{virtual void f() = 0;}? Cosa cambia, e
perche' e' \emph{undefined behavior} e non un errore di compilazione?
\vaisol{A1}
\end{domanda}

\begin{domanda}{A2}{Argomenti di default e binding}
Cosa stampa e perche'?
\begin{lstlisting}[language=C++]
struct Base { virtual void f(int x = 10) { std::cout << "Base " << x; } };
struct Der : Base { void f(int x = 20) override { std::cout << "Der " << x; } };
int main() { Base* p = new Der; p->f(); }
\end{lstlisting}
\vaisol{A2}
\end{domanda}

\begin{domanda}{A3}{Name hiding degli overload}
Perche' questo codice compila ma chiama la funzione ``sbagliata''? Correggilo
in almeno due modi diversi.
\begin{lstlisting}[language=C++]
struct B { virtual void f(int);  void f(double); };
struct D : B { void f(int) override; };
D d;  d.f(3.14);
\end{lstlisting}
\vaisol{A3}
\end{domanda}

\begin{domanda}{A4}{Dominanza ed ereditarieta' virtuale}
Cosa stampa?
\begin{lstlisting}[language=C++]
struct A { virtual void f() { std::cout << "A"; } };
struct B : virtual A { void f() override { std::cout << "B"; } };
struct C : virtual A { };
struct D : B, C { };
int main() { D d; d.f(); }
\end{lstlisting}
Cosa succederebbe \emph{senza} la parola chiave \code{virtual} nelle basi? E se
anche \code{C} facesse l'override di \code{f()}, in quale momento esatto il
programma diventa ill-formed?
\vaisol{A4}
\end{domanda}

\begin{domanda}{A5}{Ereditarieta' multipla, thunk e aggiustamento di this}
\begin{lstlisting}[language=C++]
struct A { virtual void f(); int a; };
struct B { virtual void g(); int b; };
struct C : A, B { void f() override; void g() override; };
C c;  B* pb = &c;      // (void*)pb != (void*)&c
pb->g();
\end{lstlisting}
Il \code{this} che riceve \code{C::g} deve puntare all'inizio di \code{C}, ma
\code{pb} punta al sotto-oggetto \code{B}. Chi esegue l'aggiustamento e dove e'
memorizzato? Quanti vptr contiene \code{c}? Come recuperi l'indirizzo
dell'oggetto \emph{most derived} partendo da \code{pb}?
\vaisol{A5}
\end{domanda}

\begin{domanda}{A6}{Quando una firma e' davvero un override}
Quali dei seguenti sono override legali di
\code{virtual void f(int) const noexcept;}\,? Motiva \emph{ciascuno}.
\begin{center}
\footnotesize
\begin{tabular}{|l|l|}
\hline
(a) \code{void f(int)}                    & (e) \code{void f(int) const noexcept \&\&} \\
\hline
(b) \code{void f(int) const}              & (f) \code{int f(int) const noexcept} \\
\hline
(c) \code{void f(int) const noexcept}     & (g) \code{void f(int) const noexcept override final} \\
\hline
(d) \code{void f(long) const noexcept}    & \\
\hline
\end{tabular}
\end{center}
\vaisol{A6}
\end{domanda}

\begin{domanda}{A7}{Private virtual e pattern NVI}
Spiega il pattern NVI (\emph{Non-Virtual Interface}) e di' cosa succede a
ciascuna delle tre righe finali.
\begin{lstlisting}[language=C++]
class B {
public:  void call() { f(); }
private: virtual void f() { std::cout << "B"; }
};
class D : public B {
public: void f() override { std::cout << "D"; }
};
B* p = new D;
p->call();                     // (1)
p->f();                        // (2)
static_cast<D*>(p)->f();       // (3)
\end{lstlisting}
Perche' un metodo \code{private virtual} e' comunque overridabile? Su quale
tipo viene fatto il controllo di accesso?
\vaisol{A7}
\end{domanda}

\begin{domanda}{A8}{Cinque ``perche''' sul modello a oggetti del C++}
Rispondi con motivazione tecnica (non ``perche' lo dice lo standard''):
\begin{enumerate}
    \item[(a)] perche' un \textbf{costruttore} non puo' essere \code{virtual},
          e come si implementa un \code{clone()} polimorfico con return
          covariante?
    \item[(b)] perche' una \textbf{member function template} non puo' essere
          \code{virtual}?
    \item[(c)] una funzione virtuale puo' essere \code{inline}? Con quale
          effetto reale?
    \item[(d)] cosa permette al compilatore \code{final} che \code{override}
          non permette?
    \item[(e)] perche' \code{delete[] pBase} su un array di \code{Derived} e'
          UB \emph{anche} con distruttore virtuale?
\end{enumerate}
\vaisol{A8}
\end{domanda}

\section{Ereditarieta' virtuale}

\begin{domanda}{A9}{Chi costruisce la base virtuale}
Cosa stampa? E cosa stamperebbe \code{D d;} al posto di \code{E e;}\,? E
\code{B b;}\,?
\begin{lstlisting}[language=C++]
struct A { A(const char* who) { std::cout << who << " "; } };
struct B : virtual A { B() : A("da B") {} };
struct C : virtual A { C() : A("da C") {} };
struct D : B, C  { D() : A("da D"), B(), C() {} };
struct E : D     { E() : A("da E"), D()      {} };
int main() { E e; }
\end{lstlisting}
Cosa succede se togli \code{A("da E")} dalla lista di inizializzazione di
\code{E}? E qual e' l'ordine di costruzione dei sotto-oggetti?
\vaisol{A9}
\end{domanda}

\begin{domanda}{A10}{Layout e costo della base virtuale}
Nel classico esempio \code{Animal} / \code{Mammal} / \code{WingedAnimal} /
\code{Bat} con ereditarieta' virtuale, il layout di \code{Bat} diventa
\begin{center}
\code{(vpointer, Mammal, vpointer, WingedAnimal, Bat, Animal)}
\end{center}
Rispondi: perche' il sotto-oggetto \code{Animal} finisce \textbf{in fondo} e non
all'inizio? Perche' servono \textbf{due} vpointer? Perche' l'offset fra un
sotto-oggetto \code{Mammal} e il suo \code{Animal} non e' una costante di
compilazione, come invece succede nell'ereditarieta' non virtuale? E infine:
quale delle due righe seguenti e' ill-formed, e perche'?
\begin{lstlisting}[language=C++]
Animal* pa = &bat;
Bat* p1 = static_cast<Bat*>(pa);
Bat* p2 = dynamic_cast<Bat*>(pa);
\end{lstlisting}
\vaisol{A10}
\end{domanda}

\begin{domanda}{A11}{Metodi virtuali puri nel diamante}
\begin{lstlisting}[language=C++]
struct I { virtual void f() = 0; virtual ~I() = default; };
struct L : virtual I { void f() override { std::cout << "L"; } };
struct R : virtual I { };
struct D : L, R { };
int main() { D d; d.f(); }
\end{lstlisting}
\code{D} e' istanziabile? Cosa cambia se togli \code{virtual} dalle basi di
\code{L} e \code{R}? Infine: un metodo virtuale \emph{puro} puo' avere un
\textbf{corpo}? Se si', come lo si chiama e a cosa serve? In quale caso quel
corpo e' \emph{obbligatorio}?
\vaisol{A11}
\end{domanda}

\begin{domanda}{A12}{Quando usarla, e come se la cavano gli altri tre linguaggi}
Il criterio classico e' che l'ereditarieta' virtuale serve quando l'ereditarieta'
esprime una \emph{restrizione di un insieme} e non una \emph{composizione di
parti}. Spiega il criterio, poi rispondi:
\begin{itemize}
    \item in una frase: che differenza c'e', a livello di rappresentazione in
          memoria, tra base virtuale e base non virtuale?
    \item perche' in \textbf{Java} il problema del diamante \emph{sui dati} non
          puo' nemmeno presentarsi, e quale diamante resta?
    \item perche' in \textbf{Python} si potrebbe dire che l'ereditarieta' e'
          ``sempre virtuale''?
    \item come si sceglie fra le due in \textbf{C}?
\end{itemize}
Cita infine un esempio reale di ereditarieta' virtuale nella libreria standard
del C++.
\vaisol{A12}
\end{domanda}

%==============================================================
\chapter{Java}\label{ch:java}
%==============================================================

\begin{domanda}{B1}{Campi nascosti contro metodi sovrascritti}
Cosa stampa? Ci sono tre risultati diversi: spiegali tutti e tre.
\begin{lstlisting}[language=Java]
class A { String x = "A"; String get() { return x; } String getX() { return x; } }
class B extends A { String x = "B"; String get() { return x; } }

A a = new B();
System.out.println(a.x);        // (1)
System.out.println(a.get());    // (2)
System.out.println(a.getX());   // (3)
System.out.println(((B)a).x);   // (4)
\end{lstlisting}
\vaisol{B1}
\end{domanda}

\begin{domanda}{B2}{Metodi statici e istanze}
Cosa stampa e perche'? Come si chiama tecnicamente il fenomeno (non e'
overriding)?
\begin{lstlisting}[language=Java]
class A { static String s() { return "A"; } }
class B extends A { static String s() { return "B"; } }
A a = new B();
System.out.println(a.s());
\end{lstlisting}
\vaisol{B2}
\end{domanda}

\begin{domanda}{B3}{Metodo overridabile chiamato dal costruttore}
Cosa stampa? Descrivi l'ordine \emph{esatto} di inizializzazione.
\begin{lstlisting}[language=Java]
class A { A() { init(); } void init() { } }
class B extends A {
    final int x;
    List<String> list = new ArrayList<>();
    B() { x = 5; }
    @Override void init() { System.out.println(x); System.out.println(list); }
}
new B();
\end{lstlisting}
\vaisol{B3}
\end{domanda}

\begin{domanda}{B4}{Overriding o overloading?}
Cosa stampa ciascuna riga?
\begin{lstlisting}[language=Java]
class A { void m(Object o) { System.out.println("A/Object"); } }
class B extends A { void m(String s) { System.out.println("B/String"); } }

A a = new B();  B b = new B();
a.m("hi");            // (1)
b.m("hi");            // (2)
b.m((Object)"hi");    // (3)
\end{lstlisting}
\vaisol{B4}
\end{domanda}

\begin{domanda}{B5}{Generics, erasure e bridge method}
Perche' questo lancia \code{ClassCastException} e, soprattutto, \emph{in quale
metodo} viene lanciata? Che bytecode ha generato il compilatore?
\begin{lstlisting}[language=Java]
class Node<T> { T data; void set(T t) { this.data = t; } }
class MyNode extends Node<Integer> { @Override void set(Integer i) { super.set(i); } }

Node n = new MyNode();
n.set("stringa");
\end{lstlisting}
\vaisol{B5}
\end{domanda}

\begin{domanda}{B6}{Diamante di default method}
Perche' non compila? Scrivi la correzione.
\begin{lstlisting}[language=Java]
interface A { default String f() { return "A"; } }
interface B extends A { default String f() { return "B"; } }
interface C extends A { default String f() { return "C"; } }
class D implements B, C { }
\end{lstlisting}
Poi spiega: (a) la regola \emph{class wins}; (b) perche' un'interfaccia
\textbf{non puo'} dichiarare \code{default String toString()}; (c) perche'
esiste \code{invokeinterface} distinto da \code{invokevirtual}.
\vaisol{B6}
\end{domanda}

\begin{domanda}{B7}{Le regole dell'override in Java}
Quali di queste operazioni sono legali? Motiva.
\begin{enumerate}
    \item[(a)] restringere la visibilita' \code{public} $\rightarrow$ \code{protected};
    \item[(b)] allargarla \code{protected} $\rightarrow$ \code{public};
    \item[(c)] return covariante;
    \item[(d)] aggiungere una \emph{checked exception} non dichiarata nella base;
    \item[(e)] fare l'override di un metodo \code{private} della superclasse;
    \item[(f)] un metodo \code{static} che fa l'override di uno di istanza;
    \item[(g)] scrivere \code{equals(MyType)} invece di \code{equals(Object)}:
          che bug produce con \code{HashSet.contains()}?
\end{enumerate}
\vaisol{B7}
\end{domanda}

%==============================================================
\chapter{Python}\label{ch:py}
%==============================================================

\begin{domanda}{C1}{MRO, C3 e super() cooperativo}
Dai l'MRO di \code{D} e l'output. Perche' \code{super()} dentro \code{B}
\emph{non} chiama \code{A.f}?
\begin{lstlisting}[language=Python]
class A:
    def f(self): print("A")
class B(A):
    def f(self): print("B"); super().f()
class C(A):
    def f(self): print("C"); super().f()
class D(B, C):
    def f(self): print("D"); super().f()
D().f()
\end{lstlisting}
Poi: perche' \code{class X(A, B)} con \code{class B(A)} solleva
\code{TypeError}?
\vaisol{C1}
\end{domanda}

\begin{domanda}{C2}{Lookup implicito dei metodi speciali}
Cosa stampano le due righe? Enuncia la regola.
\begin{lstlisting}[language=Python]
class A:
    def __len__(self): return 1
a = A()
a.__len__ = lambda: 99
print(len(a))
print(a.__len__())
\end{lstlisting}
Perche' questo rende difficile scrivere un \emph{proxy} trasparente basato su
\dunder{getattr}?
\vaisol{C2}
\end{domanda}

\begin{domanda}{C3}{Name mangling}
Cosa stampa? Come si chiama il meccanismo e come lo si usa per ottenere
l'equivalente di un metodo \emph{non} virtuale?
\begin{lstlisting}[language=Python]
class A:
    def __init__(self): self.__helper()
    def __helper(self): print("A helper")
class B(A):
    def __helper(self): print("B helper")
B()
\end{lstlisting}
\vaisol{C3}
\end{domanda}

\begin{domanda}{C4}{Dove vive davvero un metodo}
Analizza riga per riga.
\begin{lstlisting}[language=Python]
class A:
    def f(self): print("A.f")
class B(A):
    def f(self): print("B.f")

b = B()
b.f()                                  # (1)
A.f(b)                                 # (2)
B.f = A.f;  b.f()                      # (3)
b.f = lambda: print("inst");  b.f()    # (4)
del b.f
b.__class__ = A;  b.f()                # (5)
\end{lstlisting}
Perche' al punto (4) l'attributo di istanza vince sul metodo di classe, mentre
con una \code{property} non succederebbe?
\vaisol{C4}
\end{domanda}

\begin{domanda}{C5}{Priorita' delle operazioni riflesse}
Cosa stampa e perche'? Qual e' la regola precisa di CPython?
\begin{lstlisting}[language=Python]
class A:
    def __add__(self, other): return "A.__add__"
class B(A):
    def __radd__(self, other): return "B.__radd__"
print(A() + B())
\end{lstlisting}
\vaisol{C5}
\end{domanda}

\begin{domanda}{C6}{Il costruttore virtuale che il C++ non ha}
Cosa stampa? Perche' questo e' il modo idiomatico di scrivere factory
polimorfiche?
\begin{lstlisting}[language=Python]
class A:
    @classmethod
    def create(cls, *args): return cls(*args)
    def __repr__(self): return type(self).__name__
class B(A): pass
print(B.create())
\end{lstlisting}
Cosa cambierebbe con uno \code{@staticmethod} che restituisce \code{A()}?
\vaisol{C6}
\end{domanda}

\begin{domanda}{C7}{Metodi astratti che non astraggono}
Perche' questo \textbf{non} solleva alcun errore all'istanziazione, e come si
corregge?
\begin{lstlisting}[language=Python]
from abc import abstractmethod
class Base:
    @abstractmethod
    def f(self): ...
class Impl(Base): pass
Impl()
\end{lstlisting}
Bonus: ai fini dell'overriding, che differenza c'e' tra \dunder{init\_subclass},
una metaclasse e \code{ABCMeta.\_\_subclasshook\_\_}?
\vaisol{C7}
\end{domanda}

%==============================================================
\chapter{C}\label{ch:c}
%==============================================================

\begin{domanda}{D1}{Dispatch dinamico scritto a mano}
Implementa il dispatch dinamico in C puro: una \code{struct Shape} con vtable,
un \code{Circle} che ``eredita'', e la funzione \code{draw\_all}. Poi rispondi:
\begin{itemize}
    \item perche' il sotto-oggetto base deve essere il \textbf{primo} membro?
    \item chi inizializza il vptr, e cosa succede se te lo dimentichi
          (confronto con il C++)?
    \item come emuli un \textbf{distruttore virtuale} e un \code{dynamic\_cast}?
\end{itemize}
\vaisol{D1}
\end{domanda}

\begin{domanda}{D2}{Il cast che sembra innocuo}
Perche' questo e' formalmente scorretto in C, anche se ``funziona'' su tutti i
compilatori?
\begin{lstlisting}[language=C]
void draw_all(Shape **v, size_t n);
Circle *arr[3] = { c1, c2, c3 };
draw_all((Shape **)arr, 3);
\end{lstlisting}
Cosa dicono lo standard sulla \emph{common initial sequence} e la regola di
\emph{strict aliasing}? Perche' la garanzia ``indirizzo del primo membro =
indirizzo della struct'' non basta a salvarti qui?
\vaisol{D2}
\end{domanda}

\begin{domanda}{D3}{Traduzione letterale di virtual}
In C++ scrivi \code{virtual double area() const;}. Traduci \emph{letteralmente}
la sua rappresentazione a runtime: cosa diventa \code{this}, cosa diventa
\code{const}, cos'e' la vtable e come viene ``ereditata'' dalla classe
derivata. Qual e' il costo in byte per oggetto e per classe?
\vaisol{D3}
\end{domanda}

\begin{domanda}{D4}{Due polimorfismi diversi nella libreria standard e nel kernel}
\code{struct file\_operations} del kernel Linux e il comparatore di
\code{qsort(..., int (*cmp)(const void*, const void*))} sono due forme diverse
di polimorfismo. Quale corrisponde all'overriding e quale alla
\emph{strategy}/callback? Cosa fa la macro \code{container\_of} e a quale
operazione del C++ corrisponde?
\vaisol{D4}
\end{domanda}

\begin{domanda}{D5}{Perche' il C non puo' avere overriding automatico}
Perche' in C non puoi ottenere overriding ``automatico'' nemmeno usando
\code{static inline} negli header? Collega la risposta a: \emph{early binding},
unita' di traduzione, assenza di \emph{name mangling} e assenza di costruttori.
\vaisol{D5}
\end{domanda}

%==============================================================
\chapter{Domande trasversali}\label{ch:trasv}
%==============================================================

\begin{domanda}{E1}{Tabella comparativa del dispatch}
Compila la tabella per \textbf{C++, Java, Python, C}:
\begin{enumerate}
    \item[(i)] meccanismo di dispatch a runtime;
    \item[(ii)] costo per chiamata;
    \item[(iii)] e' possibile aggiungere o rimpiazzare un metodo a runtime?
    \item[(iv)] esiste il \emph{method hiding}?
    \item[(v)] come si chiama esplicitamente l'implementazione della base?
\end{enumerate}
\vaisol{E1}
\end{domanda}

\begin{domanda}{E2}{Virtuale per default oppure no}
In C++ le funzioni sono non virtuali per default, in Java sono virtuali per
default (salvo \code{final}, \code{private}, \code{static}), in Python tutto e'
dinamico. Argomenta il trade-off usando: \emph{zero-overhead principle},
\emph{fragile base class problem} e devirtualizzazione (CHA + inline caching
nella JVM contro LTO e \code{final} in C++).
\vaisol{E2}
\end{domanda}

\begin{domanda}{E3}{Covarianza e controvarianza}
Data \code{Base} con \code{void process(Base b)} e \code{Derived} con
\code{void process(Derived d)}: in quale dei quattro linguaggi questo e' un
override? Collega la risposta alla \textbf{covarianza/controvarianza dei
parametri} e al principio di sostituzione di Liskov. Perche' nessun linguaggio
mainstream ammette parametri controvarianti?
\vaisol{E3}
\end{domanda}

\begin{domanda}{E4}{Sette termini da non confondere}
Definisci con precisione e distingui: \emph{overriding}, \emph{overloading},
\emph{hiding/shadowing}, \emph{dominanza}, \emph{bridge method}, \emph{thunk},
\emph{final overrider}. Per ciascuno indica in quale dei quattro linguaggi
esiste.
\vaisol{E4}
\end{domanda}

%==============================================================
\chapter{Soluzioni}\label{ch:sol}
%==============================================================

\section{Soluzioni C++}

\begin{soluzione}{A1 --- Chiamate virtuali in costruttore e distruttore}
\label{sol:A1}
Output: \code{A() A::f A::g}\,.

Durante il corpo del costruttore di \code{A} il tipo dinamico dell'oggetto
\textbf{e'} \code{A}: il vptr e' stato impostato a quello di \code{A} e verra'
riscritto solo all'inizio del costruttore di \code{B}. Quindi \code{f()}
risolve ad \code{A::f} --- e per fortuna, perche' \code{B::x} non e' ancora
inizializzato e \code{B::f} lo leggerebbe.

Simmetricamente \code{\~{}A()} gira \emph{dopo} \code{\~{}B()}: quando si entra
nel corpo di \code{\~{}A} il vptr e' gia' stato riportato ad \code{A}, quindi
\code{g()} chiama \code{A::g}.

\textbf{Regola:} mai chiamare funzioni virtuali da costruttore o distruttore
aspettandosi dispatch dinamico. Il tipo dinamico ``cresce'' durante la
costruzione e ``decresce'' durante la distruzione.

Con \code{f()} pura, la chiamata dal costruttore di \code{A} e'
\textbf{undefined behavior} (\emph{[class.abstract]}). Non e' un errore di
compilazione perche' la chiamata e' sintatticamente lecita e in generale il
compilatore non puo' dimostrare che verra' eseguita (potrebbe essere dentro un
\code{if}, o in una funzione chiamata dal costruttore). In pratica quasi tutte
le implementazioni mettono nella vtable di \code{A} il simbolo
\code{\_\_cxa\_pure\_virtual}, che termina il programma con il messaggio
``pure virtual method called''.
\end{soluzione}

\begin{soluzione}{A2 --- Argomenti di default e binding}
\label{sol:A2}
Stampa \code{Der 10}.

Il \emph{corpo} della funzione e' scelto dinamicamente (tramite vtable), ma gli
\textbf{argomenti di default sono sostituiti a tempo di compilazione} in base
al \emph{tipo statico} dell'espressione, che qui e' \code{Base*}. Il
compilatore genera di fatto \code{p->f(10)} e poi il dispatch porta a
\code{Der::f}.

Da qui la regola pratica: \textbf{mai ridefinire gli argomenti di default in un
override}. Se servono default diversi, si usa NVI (vedi A7): la funzione
pubblica non virtuale porta il default, quella virtuale privata non ne ha.
\end{soluzione}

\begin{soluzione}{A3 --- Name hiding degli overload}
\label{sol:A3}
\code{D::f(int)} \textbf{nasconde} tutti gli overload \code{B::f}: il
\emph{name lookup} risale gli scope e si ferma al \emph{primo} che contiene il
nome \code{f}, cioe' \code{D}. Solo dopo, sull'insieme trovato, parte
l'\emph{overload resolution}. Quindi l'unico candidato e' \code{D::f(int)}, e
\code{3.14} viene convertito a \code{int} (con troncamento silenzioso).

Correzioni possibili:
\begin{lstlisting}[language=C++]
struct D : B {
    using B::f;              // (1) reintroduce tutti gli overload della base
    void f(int) override;
};
d.B::f(3.14);                // (2) chiamata qualificata
// (3) ridichiarare esplicitamente void f(double) in D
\end{lstlisting}
Il flag \code{-Woverloaded-virtual} di GCC/Clang intercetta il caso.
\end{soluzione}

\begin{soluzione}{A4 --- Dominanza ed ereditarieta' virtuale}
\label{sol:A4}
Stampa \code{B}.

Con basi virtuali esiste \textbf{un solo} sotto-oggetto \code{A} condiviso, e
la regola di \textbf{dominanza} stabilisce che \code{B::f} domina \code{A::f}
(perche' \code{B} deriva da \code{A}). Il \emph{final overrider} di \code{A::f}
in \code{D} e' quindi unico: nessuna ambiguita'.

Senza \code{virtual} nelle basi ci sarebbero \textbf{due} sotto-oggetti
\code{A} distinti, quindi due \code{A::f} indipendenti: \code{d.f()} e'
ambiguo (errore di compilazione) e servirebbe \code{d.B::f()}.

Se anche \code{C} facesse l'override, non esisterebbe piu' un final overrider
unico: il programma e' \textbf{ill-formed alla definizione di \code{D}} (non al
punto di chiamata), a meno che \code{D} non fornisca il proprio override, che
allora domina entrambi.
\end{soluzione}

\begin{soluzione}{A5 --- Ereditarieta' multipla, thunk e aggiustamento di this}
\label{sol:A5}
L'aggiustamento del \code{this} lo esegue un \textbf{thunk}: nella
\emph{vtable secondaria} (quella associata al sotto-oggetto \code{B} dentro
\code{C}) la entry per \code{g} non punta a \code{C::g}, ma a un piccolo
trampolino che fa \code{this -= offsetof(C, B-subobject)} e poi salta a
\code{C::g}. Nell'Itanium ABI si chiama \emph{non-virtual this adjustment};
altre ABI memorizzano l'offset accanto al puntatore a funzione.

\code{c} contiene \textbf{due vptr}: uno per il sotto-oggetto primario
(\code{A}) e uno per \code{B}. Ecco perche' \code{sizeof(C)} e' maggiore di
quanto ci si aspetterebbe dai soli campi.

L'indirizzo dell'oggetto \emph{most derived} si ottiene con
\code{dynamic\_cast<void*>(pb)}, che usa l'\emph{offset-to-top} memorizzato
nella vtable prima delle entry dei metodi.
\end{soluzione}

\begin{soluzione}{A6 --- Quando una firma e' davvero un override}
\label{sol:A6}
\begin{itemize}
    \item[(a)] \code{void f(int)} --- \textbf{no}: perde \code{const}, quindi
          la firma e' diversa. Diventa un overload che \emph{nasconde} il
          virtuale; con \code{override} il compilatore da' errore.
    \item[(b)] \code{void f(int) const} --- \textbf{no}: da C++17 \code{noexcept}
          fa parte del \emph{tipo} della funzione, e un override non puo' essere
          \emph{meno} restrittivo della base. (Il contrario --- base non
          \code{noexcept}, derivata \code{noexcept} --- e' invece legale.)
    \item[(c)] \code{void f(int) const noexcept} --- \textbf{si'}, e' l'override
          corretto.
    \item[(d)] \code{void f(long) const noexcept} --- \textbf{no}: parametri
          diversi, quindi overload + hiding.
    \item[(e)] con \code{\&\&} --- \textbf{no}: i \emph{ref-qualifier} fanno
          parte della firma e devono coincidere esattamente.
    \item[(f)] \code{int f(int) const noexcept} --- \textbf{no}: il return type
          e' diverso e non covariante. La covarianza vale solo per puntatori o
          riferimenti a classi in relazione di derivazione (e non per
          \code{int} contro \code{void}).
    \item[(g)] \code{override final} in una classe gia' \code{final} ---
          legale, ma ridondante.
\end{itemize}
\end{soluzione}

\begin{soluzione}{A7 --- Private virtual e pattern NVI}
\label{sol:A7}
(1) stampa \code{D}; \quad (2) \textbf{errore di compilazione}; \quad
(3) compila e stampa \code{D}.

Il \emph{name lookup} e il \emph{controllo di accesso} avvengono sul
\textbf{tipo statico} e sono completamente ortogonali al dispatch dinamico, che
avviene sul tipo dinamico. Per questo un \code{private virtual} e'
perfettamente overridabile: la classe derivata non ha bisogno di
\emph{chiamarlo}, solo di \emph{ridefinirlo}.

\textbf{NVI:} l'interfaccia pubblica e' non virtuale (\code{call()}) e i punti
di variazione sono virtuali privati. In questo modo la base mantiene il
controllo su pre/post-condizioni, logging, lock, invarianti --- ed e' libera di
cambiare il numero e la forma dei punti di estensione senza rompere le
derivate. Corollario noto (Sutter): in una classe base pubblica, l'unica
funzione che dovrebbe essere \code{public} \emph{e} \code{virtual} e' il
distruttore.
\end{soluzione}

\begin{soluzione}{A8 --- Cinque ``perche''' sul modello a oggetti del C++}
\label{sol:A8}
\textbf{(a)} Il dispatch virtuale richiede un oggetto gia' costruito con un
vptr valido, ed e' esattamente il costruttore a installarlo: un costruttore
virtuale sarebbe circolare. L'idioma sostitutivo e' il \emph{virtual
constructor} di Coplien:
\begin{lstlisting}[language=C++]
struct Base {
    virtual Base* clone() const { return new Base(*this); }
    virtual ~Base() = default;
};
struct Der : Base {
    Der* clone() const override { return new Der(*this); }  // return covariante
};
\end{lstlisting}

\textbf{(b)} Una member template genera un numero potenzialmente illimitato di
istanziazioni, scoperte in unita' di traduzione diverse e anche molto dopo la
definizione della classe: la vtable dovrebbe avere dimensione e layout non
determinabili nel punto in cui la classe viene definita. (Una \emph{class}
template puo' invece avere tutte le funzioni virtuali che vuole: per ogni
istanziazione concreta l'insieme dei metodi e' noto.)

\textbf{(c)} Si'. \code{inline} resta utile per le chiamate risolte
staticamente (\code{obj.f()} su un oggetto, \code{Base::f()} qualificato, o
dopo devirtualizzazione); il compilatore deve comunque emettere anche una copia
\emph{out-of-line}, perche' l'indirizzo della funzione finisce nella vtable.

\textbf{(d)} \code{override} e' solo un \emph{controllo}: ``verifica che io
stia davvero sovrascrivendo qualcosa''. \code{final} e' invece
un'\emph{informazione semantica} --- nessuno potra' piu' fare l'override ---
che permette al compilatore di \textbf{devirtualizzare} la chiamata e quindi di
inlinearla, anche senza LTO o analisi dell'intero programma.

\textbf{(e)} L'aritmetica dei puntatori negli array usa \code{sizeof} del
\textbf{tipo statico}: \code{pBase[1]} verrebbe calcolato con passo
\code{sizeof(Base)} invece di \code{sizeof(Derived)}. Il distruttore virtuale
non c'entra nulla: il problema e' l'indicizzazione, non la distruzione. La
soluzione e' \code{std::vector<std::unique\_ptr<Base>{}>}.
\end{soluzione}

\begin{soluzione}{A9 --- Chi costruisce la base virtuale}
\label{sol:A9}
Stampa \code{da E}, \textbf{una volta sola}.

La regola e' che \textbf{la base virtuale viene inizializzata dal costruttore
della classe \emph{most derived}}, cioe' dalla classe dell'oggetto che si sta
effettivamente creando. Tutti gli altri mem-initializer per \code{A} ---
quelli in \code{B}, \code{C} e \code{D} --- vengono semplicemente
\emph{ignorati}, perche' in questa costruzione quelle classi sono sotto-oggetti,
non l'oggetto piu' derivato.

Quindi: \code{D d;} stampa \code{da D}, e \code{B b;} stampa \code{da B}. Lo
stesso mem-initializer \code{A("da B")} viene usato o ignorato a seconda del
ruolo che \code{B} ha in quella particolare costruzione: ecco perche' scriverlo
in \code{B} non e' inutile, serve per il caso in cui \code{B} sia standalone.

Se togli \code{A("da E")} da \code{E}, il compilatore deve inizializzare
\code{A} per default --- ma \code{A} non ha costruttore di default, quindi il
programma e' \textbf{ill-formed}. Conseguenza pratica: una base virtuale
conviene che abbia un costruttore di default, altrimenti \emph{ogni} classe
concreta della gerarchia (a qualunque profondita') e' obbligata a nominarla
esplicitamente. E' uno dei costi ``sociali'' dell'ereditarieta' virtuale.

\textbf{Ordine di costruzione:} prima tutte le \emph{basi virtuali}, in ordine
di visita depth-first da sinistra a destra della gerarchia della classe most
derived; poi le basi non virtuali nell'ordine di dichiarazione; poi i membri
nell'ordine di dichiarazione; infine il corpo del costruttore. Qui:
\code{A}, \code{B}, \code{C}, \code{D}, \code{E}. La distruzione avviene in
ordine esattamente inverso.
\end{soluzione}

\begin{soluzione}{A10 --- Layout e costo della base virtuale}
\label{sol:A10}
\code{Animal} finisce \textbf{in fondo} proprio perche' e' \emph{condiviso}: i
sotto-oggetti \code{Mammal} e \code{WingedAnimal} devono conservare il layout
che avrebbero da soli (cosi' il codice compilato per \code{Mammal} continua a
funzionare), quindi l'unica copia della base comune viene messa in una posizione
che non appartiene a nessuno dei due rami. Nel caso non virtuale, invece, ogni
ramo si porta dietro la \emph{propria} copia della base, incorporata al proprio
inizio.

Servono \textbf{due vpointer} perche' ciascun ramo che eredita virtualmente ha
bisogno di sapere ``quanto dista da qui il mio \code{Animal}'', e questa
distanza e' diversa a seconda di come il sotto-oggetto e' immerso nell'oggetto
completo. L'informazione si chiama \textbf{vbase offset} e sta nella vtable,
accanto ai puntatori ai metodi.

Perche' non e' una costante di compilazione? Perche' quando il compilatore
traduce un metodo di \code{Mammal} non sa se quel \code{Mammal} sara' un oggetto
autonomo (e allora \code{Animal} sta subito dopo il vptr) oppure il sotto-oggetto
di un \code{Bat} (e allora \code{Animal} sta molto piu' in la', dopo
\code{WingedAnimal} e dopo i membri di \code{Bat}). Con ereditarieta' non
virtuale l'offset base/derivata e' invece fissato una volta per tutte a compile
time --- ed e' esattamente la differenza di costo fra i due meccanismi: nel caso
virtuale l'accesso a un membro della base costa un'indirezione in piu'. Nel
conto di Wikipedia, la dimensione dell'oggetto \code{Bat} cresce di due
puntatori rispetto alla versione non virtuale, in cambio dell'unicita' di
\code{Animal} e dell'assenza di ambiguita'.

\textbf{Il cast:} e' ill-formed la riga con \code{static\_cast}. Un
\code{static\_cast} da base virtuale a derivata e' \emph{vietato dallo
standard}, e per il motivo appena detto: il compilatore dovrebbe sottrarre un
offset che non conosce staticamente. Serve \code{dynamic\_cast}, che legge
l'informazione dalla vtable a runtime (ed e' anche l'unico caso in cui
\code{dynamic\_cast} e' obbligatorio pur non essendoci alcuna incertezza sul
tipo). Per lo stesso motivo non si puo' usare \code{offsetof} su un membro della
base virtuale, e i puntatori a membro della base virtuale hanno una
rappresentazione piu' complessa.
\end{soluzione}

\begin{soluzione}{A11 --- Metodi virtuali puri nel diamante}
\label{sol:A11}
\textbf{Si', \code{D} e' istanziabile} e stampa \code{L}. Con le basi virtuali
esiste un solo sotto-oggetto \code{I}, e per la regola di dominanza (A4) il suo
\emph{final overrider} unico e' \code{L::f}. Il ramo \code{R} non ha bisogno di
implementare nulla: eredita l'implementazione ``dall'altra parte'' del diamante.
E' esattamente il punto sottolineato da Wikipedia --- con ereditarieta' virtuale
il metodo puro va implementato \emph{una volta sola} nella gerarchia.

Togliendo \code{virtual} dalle basi, ci sono \textbf{due} sotto-oggetti \code{I}
indipendenti: quello raggiunto via \code{R} ha ancora \code{f} pura, quindi
\code{D} resta \textbf{astratta} e \code{D d;} non compila. E anche se
implementassi \code{f} in \code{R}, la chiamata \code{d.f()} sarebbe ambigua e
servirebbe \code{d.L::f()}.

\textbf{Un metodo puro puo' avere un corpo.} \code{= 0} dichiara solo che le
derivate concrete devono fornire un override; non vieta di dare
un'implementazione di \emph{default}:
\begin{lstlisting}[language=C++]
struct I { virtual void f() = 0; };
void I::f() { std::cout << "comportamento di base"; }
struct D : I { void f() override { I::f(); std::cout << " + extra"; } };
\end{lstlisting}
Il corpo si puo' invocare solo con \textbf{chiamata qualificata}
(\code{I::f()}), mai per dispatch dinamico: la entry in vtable resta quella
``pura''. Serve come fallback riusabile dalle derivate --- ed e' lo stesso
schema di Python, dove un \code{@abstractmethod} puo' avere un corpo che la
sottoclasse richiama con \code{super()}.

Il corpo e' \textbf{obbligatorio} in un caso: il \emph{distruttore virtuale
puro}. \code{virtual \~{}I() = 0;} rende la classe astratta (idioma utile quando
non ci sono altri metodi da rendere puri), ma il distruttore di ogni derivata
chiama comunque quello della base, quindi la definizione deve esistere:
\code{I::\~{}I() \{\}}.
\end{soluzione}

\begin{soluzione}{A12 --- Quando usarla, e come se la cavano gli altri tre linguaggi}
\label{sol:A12}
\textbf{Il criterio.} Se l'ereditarieta' esprime una \emph{restrizione di un
insieme} --- un \code{Bat} \emph{e'} un \code{Animal}, e ne esiste uno solo per
ogni pipistrello --- allora la base va condivisa, quindi \code{virtual}. Se
invece l'ereditarieta' viene usata per \emph{comporre parti} indipendenti (un
\code{Bat} ``contiene'' le caratteristiche del volo e quelle dei mammiferi come
due moduli distinti), avere due copie e' coerente, e la base non virtuale va
bene. In pratica: la duplicazione e' un bug quando la base rappresenta
\emph{identita'}, e' accettabile quando rappresenta \emph{implementazione}.

\textbf{Rappresentazione in memoria, in una frase:} una base non virtuale e'
\emph{incorporata} (una copia per ogni ramo, offset costante), una base virtuale
e' \emph{condivisa} (una copia sola, raggiunta tramite un offset letto a runtime)
--- cioe' concettualmente ``per valore'' contro ``per riferimento''.

\textbf{Java.} Una classe puo' estendere una sola classe, quindi lo stato non
puo' arrivare da due strade: il diamante \emph{sui dati} e' impedito per
costruzione. Resta il diamante \emph{sul comportamento}, quello dei
\code{default method} delle interfacce (B6), che il linguaggio non risolve
automaticamente ma obbliga a risolvere a mano con un override esplicito e
\code{X.super.f()}. Le interfacce sono in un certo senso ``sempre virtuali'':
non portando stato, non c'e' niente da duplicare.

\textbf{Python.} L'ereditarieta' multipla di stato e' permessa, ma gli attributi
vivono in \emph{un solo} \code{\_\_dict\_\_} per istanza, e la linearizzazione C3
mette ogni classe base \emph{una volta sola} nell'MRO. Quindi Python si comporta
di fatto come se ogni base fosse virtuale: c'e' una sola \code{A}, e se tutti
cooperano con \code{super()} il suo \dunder{init} viene eseguito esattamente una
volta (C1). Il prezzo lo si paga altrove: l'ordine di inizializzazione dipende
dall'MRO della classe concreta, esattamente come in C++ la base virtuale e'
costruita dalla most derived (A9).

\textbf{C.} La scelta e' esplicita e si vede: base non virtuale $=$
\code{struct A base;} incorporata come primo membro (D1); base virtuale $=$
\code{struct A *base;}, un puntatore alla copia condivisa. Il costo
dell'ereditarieta' virtuale in C++ e' letteralmente il costo di quella
sostituzione --- un'indirezione in piu' e un puntatore in piu' per ramo.

\textbf{Esempio reale nella libreria standard:} la gerarchia degli stream.
\code{basic\_istream} e \code{basic\_ostream} derivano \emph{virtualmente} da
\code{basic\_ios}, cosicche' \code{basic\_iostream} (da cui \code{std::stringstream},
\code{std::fstream}) abbia un solo stato di stream --- un solo buffer, un solo
insieme di flag di errore --- invece di due. Se l'ereditarieta' non fosse
virtuale, \code{ss.fail()} potrebbe dare risposte diverse a seconda del ramo da
cui la si interroga.
\end{soluzione}

\section{Soluzioni Java}

\begin{soluzione}{B1 --- Campi nascosti contro metodi sovrascritti}
\label{sol:B1}
\begin{enumerate}
    \item[(1)] \code{A} --- i \textbf{campi non sono polimorfici}: l'accesso e'
          risolto a tempo di compilazione sul tipo statico della variabile
          (\code{A a}), con l'istruzione \code{getfield A.x}.
    \item[(2)] \code{B} --- i metodi si': \code{invokevirtual} risolve sul tipo
          dinamico.
    \item[(3)] \code{A} --- \code{getX()} e' definito \emph{dentro} \code{A}, e
          nel suo corpo l'identificatore \code{x} si riferisce ad \code{A.x}.
    \item[(4)] \code{B}.
\end{enumerate}
Il punto chiave e' che \code{B.x} non \emph{sostituisce} \code{A.x}: lo
\textbf{nasconde}. Nell'oggetto esistono \emph{due} campi distinti, e \code{A.x}
resta raggiungibile da dentro \code{B} con \code{super.x} oppure
\code{((A)this).x}. Morale: i campi si dichiarano \code{private} e si espongono
con metodi, altrimenti l'ereditarieta' diventa illeggibile.
\end{soluzione}

\begin{soluzione}{B2 --- Metodi statici e istanze}
\label{sol:B2}
Stampa \code{A}. E' \textbf{hiding di metodo statico}, non overriding: la
chiamata viene compilata in \code{invokestatic A.s()} sulla base del tipo
statico di \code{a}; l'oggetto non viene nemmeno usato (\code{javac} emette il
warning ``static method should be qualified by type name'').

I metodi \code{static} sono legati staticamente \emph{per definizione}: non
esiste un ricevitore su cui fare dispatch, quindi non esiste nemmeno un slot in
vtable. Lo stesso vale per \code{private} (\code{invokespecial}) e per i metodi
\code{final}, che pero' uno slot ce l'hanno.
\end{soluzione}

\begin{soluzione}{B3 --- Metodo overridabile chiamato dal costruttore}
\label{sol:B3}
Stampa \code{0} e poi \code{null}.

Ordine di inizializzazione:
\begin{enumerate}
    \item \code{super()} --- e qui sta il punto: il costruttore di \code{A}
          chiama \code{init()}, che e' \textbf{gia'} dispatchato dinamicamente
          su \code{B}. A differenza del C++ (A1), in Java il tipo dinamico
          dell'oggetto e' quello finale \emph{fin dal primo istante}: non
          esiste il ``vptr che cresce'';
    \item \emph{dopo} il ritorno da \code{super()} vengono eseguiti gli
          inizializzatori di campo e i blocchi di istanza di \code{B}, in ordine
          testuale;
    \item infine il corpo del costruttore di \code{B}.
\end{enumerate}
Percio' durante \code{init()} i campi di \code{B} valgono ancora il default:
\code{x == 0} \emph{anche se dichiarato \code{final}}, e \code{list == null}
(NullPointerException se ci si fa qualcosa).

E' l'errore classico documentato in \emph{Effective Java} (item 19):
\textbf{mai invocare metodi overridabili dal costruttore}. Se serve, si usa un
metodo \code{private} o \code{final}, oppure un factory method che costruisce e
\emph{poi} inizializza.
\end{soluzione}

\begin{soluzione}{B4 --- Overriding o overloading?}
\label{sol:B4}
(1) \code{A/Object} \quad (2) \code{B/String} \quad (3) \code{A/Object}.

\code{B.m(String)} \textbf{non} fa l'override di \code{A.m(Object)}: le firme
sono diverse, quindi e' un \emph{overload} che si aggiunge al metodo ereditato.

La risoluzione dell'overload avviene a \textbf{tempo di compilazione} sul tipo
statico: con \code{A a} l'unico candidato visibile e' \code{m(Object)}; con
\code{B b} i candidati sono due e vince il piu' specifico, \code{m(String)}; col
cast esplicito il tipo statico dell'argomento torna \code{Object}.

\textbf{Slogan da ripetere all'esame:} \emph{overload = statico, override =
dinamico}. Mescolare i due e' la fonte di bug piu' comune sull'argomento (vedi
anche B7g con \code{equals}).
\end{soluzione}

\begin{soluzione}{B5 --- Generics, erasure e bridge method}
\label{sol:B5}
L'eccezione viene lanciata \textbf{dentro il bridge method sintetico
\code{MyNode.set(Object)}}, non nel codice scritto dal programmatore.

Per \emph{erasure}, \code{Node.set(T)} diventa \code{set(Object)}. Ma
\code{MyNode.set(Integer)} ha un descriptor diverso, quindi da solo
\emph{non} sarebbe un override per la JVM. Il compilatore genera allora in
\code{MyNode}:
\begin{lstlisting}[language=Java]
// ACC_BRIDGE | ACC_SYNTHETIC
void set(Object o) { this.set((Integer) o); }
\end{lstlisting}
Poiche' \code{n} e' dichiarata di tipo \textbf{raw}, il compilatore non inserisce
alcun \code{checkcast} al call site: viene emessa
\code{invokevirtual Node.set(Object)}, il dispatch arriva al bridge, il bridge
esegue \code{checkcast Integer} su una \code{String} e lancia
\code{ClassCastException}. Il fenomeno si chiama \emph{heap pollution}, ed e' il
prezzo della retrocompatibilita' scelta da Java (erasure) invece della
reificazione dei generici (come in C\#).

I bridge method vengono generati anche per i \textbf{return covarianti}: se la
base restituisce \code{Object} e la derivata \code{String}, servono due metodi.
\end{soluzione}

\begin{soluzione}{B6 --- Diamante di default method}
\label{sol:B6}
\code{D} eredita due default method \code{f()} egualmente specifici e non
correlati fra loro: \code{B} e \code{C} sono entrambe sotto-interfacce di
\code{A}, ma nessuna delle due e' sotto-interfaccia dell'altra, quindi non c'e'
un criterio per preferirne una. Errore: \emph{inherits unrelated defaults}.
Correzione:
\begin{lstlisting}[language=Java]
class D implements B, C {
    @Override public String f() { return B.super.f(); }
}
\end{lstlisting}

\textbf{(a) Class wins.} Un metodo concreto ereditato da una \emph{superclasse}
batte sempre un default method, anche se l'interfaccia sembra ``piu' vicina''.
La regola esiste per garantire la retrocompatibilita': aggiungere un default a
un'interfaccia gia' pubblicata non deve mai cambiare il comportamento del codice
esistente.

\textbf{(b)} I metodi di \code{Object} sono ereditati da \emph{ogni} classe:
per la regola ``class wins'' un default \code{toString()} sarebbe sempre
scavalcato e quindi inutile o ambiguo. Il JLS lo vieta esplicitamente a tempo di
compilazione.

\textbf{(c)} Con l'ereditarieta' singola delle classi, ogni metodo virtuale ha
uno \textbf{slot a indice fisso} nella vtable, identico in tutte le
sottoclassi: \code{invokevirtual} e' quindi un accesso a offset costante. Le
interfacce invece possono essere implementate da classi con gerarchie
completamente diverse, e non esiste un indice comune: \code{invokeinterface}
richiede una ricerca (itable, itable stub, o inline cache costruita dal JIT).
\end{soluzione}

\begin{soluzione}{B7 --- Le regole dell'override in Java}
\label{sol:B7}
\begin{itemize}
    \item[(a)] \textbf{illegale}: non si puo' restringere la visibilita', o si
          violerebbe la sostituibilita' (un client che vede \code{Base} potrebbe
          chiamare il metodo).
    \item[(b)] \textbf{legale}: allargare si puo'.
    \item[(c)] \textbf{legale} da Java 5 (implementata con bridge method, B5).
    \item[(d)] \textbf{illegale} per le \emph{checked} exception: si possono
          solo restringere o omettere. Le \emph{unchecked} sono libere, perche'
          non fanno parte del contratto verificato dal compilatore.
    \item[(e)] \textbf{impossibile}: i metodi \code{private} non sono ereditati
          ne' virtuali (\code{invokespecial}); definirne uno omonimo nella
          derivata crea un metodo del tutto indipendente.
    \item[(f)] \textbf{illegale} in entrambe le direzioni: \code{static} non
          puo' nascondere un metodo di istanza ne' viceversa.
    \item[(g)] \code{equals(MyType)} e' un \emph{overload}. Tutte le collezioni
          (\code{HashSet}, \code{HashMap}, \code{List.contains}) chiamano
          \code{equals(Object)}, che rimane quello di \code{Object}, cioe'
          l'identita' di riferimento: \code{contains} restituisce \code{false}
          su oggetti logicamente uguali. Annotare sempre con \code{@Override}
          --- sulla firma sbagliata da' errore di compilazione --- e ricordare
          che \code{equals} va sempre accompagnato da \code{hashCode}.
\end{itemize}
\end{soluzione}

\section{Soluzioni Python}

\begin{soluzione}{C1 --- MRO, C3 e super() cooperativo}
\label{sol:C1}
MRO: \code{D $\rightarrow$ B $\rightarrow$ C $\rightarrow$ A $\rightarrow$
object} (linearizzazione C3). Output: \code{D B C A}.

Il punto cruciale e' che \textbf{\code{super()} non significa ``la
superclasse''}, ma ``il prossimo nell'MRO \emph{dell'oggetto}''. Dentro
\code{B.f}, con \code{type(self) is D}, il successivo dopo \code{B} nella lista
e' \code{C}, non \code{A}. E' esattamente questo che rende possibile
l'ereditarieta' multipla \emph{cooperativa}: ogni classe chiama \code{super()}
una volta sola e ogni implementazione di \code{f} viene eseguita esattamente una
volta, in ordine deterministico.

Corollario pratico: in una gerarchia cooperativa tutti i metodi devono avere
firme compatibili (tipicamente \code{*args, **kwargs}), altrimenti la catena si
rompe a meta'.

\code{class X(A, B)} con \code{class B(A)}: C3 deve rispettare sia l'ordine
locale delle basi (\code{A} prima di \code{B}) sia la monotonicita' (\code{B}
prima di \code{A}, perche' \code{B} deriva da \code{A}). I due vincoli sono
contraddittori, quindi si ottiene
\code{TypeError: Cannot create a consistent method resolution order (MRO)}.
\end{soluzione}

\begin{soluzione}{C2 --- Lookup implicito dei metodi speciali}
\label{sol:C2}
\code{len(a)} stampa \textbf{1}; \code{a.\_\_len\_\_()} stampa \textbf{99}.

L'invocazione \textbf{implicita} dei metodi speciali salta il dizionario
dell'istanza e cerca lo slot direttamente sul \textbf{tipo}
(\code{type(a).\_\_len\_\_}). E' cio' che permette a CPython di usare gli slot
scritti in C (\code{tp\_as\_sequence->sq\_length}) e di non pagare una ricerca
per attributo a ogni operazione. La ricerca \emph{esplicita} per attributo
segue invece la regola normale: istanza, poi tipo, poi MRO.

Conseguenza per i proxy: \dunder{getattr} non viene \emph{mai} consultato per le
operazioni implicite (\code{len(p)}, \code{p + q}, \code{with p:},
\code{p[0]}, \code{iter(p)}\dots), perche' la ricerca non passa dal
\dunder{getattribute} dell'istanza. Per questo \code{unittest.mock.MagicMock}
deve pre-definire tutti i dunder \emph{sul tipo}, e i proxy davvero trasparenti
si costruiscono generando dinamicamente una classe.
\end{soluzione}

\begin{soluzione}{C3 --- Name mangling}
\label{sol:C3}
Stampa \code{A helper}.

E' il \textbf{name mangling}: dentro il corpo di \code{A}, il compilatore
riscrive \code{self.\_\_helper} in \code{self.\_A\_\_helper}; \code{B} definisce
invece \code{\_B\_\_helper}. I due nomi non collidono, quindi non c'e' nessun
override.

E' l'unico modo idiomatico in Python per ottenere un metodo
\textbf{effettivamente non virtuale}: serve per implementare NVI, e soprattutto
per evitare collisioni accidentali di nomi nei \emph{mixin} e nelle gerarchie
profonde. Non e' un meccanismo di sicurezza: \code{b.\_A\_\_helper()} funziona
benissimo. Con un solo underscore (\code{\_helper}) non c'e' mangling: e' solo
una convenzione, e l'override avviene normalmente.
\end{soluzione}

\begin{soluzione}{C4 --- Dove vive davvero un metodo}
\label{sol:C4}
\begin{enumerate}
    \item[(1)] \code{B.f} --- dispatch normale lungo \code{type(b).\_\_mro\_\_}.
    \item[(2)] \code{A.f} --- chiamata esplicita dell'implementazione della base,
          equivalente a \code{super(B, b).f()}. In Python 3 \code{A.f} e' una
          funzione semplice, quindi \code{self} va passato a mano.
    \item[(3)] \code{A.f} --- le classi sono oggetti mutabili: si e' sostituito
          il metodo a runtime (\emph{monkey patching}). Nessun altro linguaggio
          del corso lo permette senza strumenti esterni.
    \item[(4)] \code{inst} --- l'attributo di istanza vince (e la lambda non
          riceve \code{self}: e' una funzione qualunque, non un metodo).
    \item[(5)] \code{A.f} --- riassegnare \dunder{class} cambia il tipo dinamico
          dell'oggetto e quindi il dispatch.
\end{enumerate}
Il motivo del punto (4) e' il \textbf{protocollo dei descrittori}: le funzioni
sono descrittori \emph{non-data} (definiscono solo \dunder{get}), e per questi
il \code{\_\_dict\_\_} dell'istanza ha la precedenza. Una \code{property} e'
invece un descrittore \emph{data} (definisce \dunder{get} \emph{e}
\dunder{set}), e i descrittori data vincono \emph{sempre} sul dizionario
dell'istanza: e' proprio cosi' che una property riesce a intercettare le
assegnazioni.
\end{soluzione}

\begin{soluzione}{C5 --- Priorita' delle operazioni riflesse}
\label{sol:C5}
Stampa \code{B.\_\_radd\_\_}.

Regola generale: per \code{x + y} Python prova \code{type(x).\_\_add\_\_} e poi
\code{type(y).\_\_radd\_\_}. \textbf{Ma} se \code{type(y)} e' una
\emph{sottoclasse propria} di \code{type(x)} \emph{e} fornisce un
\code{\_\_radd\_\_} diverso da quello (eventuale) di \code{type(x)}, il metodo
riflesso viene provato \textbf{per primo}. In CPython il controllo e' la
funzione \code{method\_is\_overloaded}: se il tipo di sinistra non ha affatto
\code{\_\_radd\_\_}, quello della sottoclasse conta come ``overloaded'' e vince.

La motivazione di design e' che una sottoclasse deve sempre poter avere
l'ultima parola sulle operazioni che la coinvolgono, altrimenti non potrebbe
raffinare il comportamento ereditato (si pensi a \code{int} e a una sottoclasse
che rappresenta unita' di misura).
\end{soluzione}

\begin{soluzione}{C6 --- Il costruttore virtuale che il C++ non ha}
\label{sol:C6}
Stampa \code{B}.

In un \code{classmethod}, \code{cls} e' legato al tipo su cui e' avvenuta
\emph{effettivamente} la chiamata (\code{B}), quindi \code{cls()} costruisce una
\code{B}: i classmethod sono \textbf{polimorfici}. E' letteralmente il
\emph{costruttore virtuale} che in C++ va emulato con \code{clone()} (A8a) e
che in Java richiede una factory esplicita o la reflection.

Con uno \code{@staticmethod} che fa \code{return A()} si perde l'informazione
sul tipo: la factory restituirebbe sempre una \code{A} e ogni sottoclasse
sarebbe costretta a riscriverla. Vale anche l'inverso: dentro un classmethod,
\code{super().create()} funziona e propaga correttamente \code{cls}.
\end{soluzione}

\begin{soluzione}{C7 --- Metodi astratti che non astraggono}
\label{sol:C7}
\code{Impl()} non solleva nulla perche' il controllo dei metodi astratti non e'
fatto dal decoratore, ma da \code{ABCMeta} in fase di creazione della classe:
e' la metaclasse a popolare l'insieme \code{\_\_abstractmethods\_\_}, che
\code{object.\_\_new\_\_} poi consulta. Il decoratore \code{@abstractmethod} da
solo si limita a impostare \code{\_\_isabstractmethod\_\_ = True} sulla
funzione, e senza \code{ABCMeta} nessuno legge quel flag.

Correzione: \code{class Base(abc.ABC):} oppure
\code{class Base(metaclass=abc.ABCMeta):}.

\textbf{Bonus.} \dunder{init\_subclass} e' un hook \emph{per sottoclasse},
invocato alla creazione di ogni classe figlia: e' il modo leggero per
validare o registrare le sottoclassi senza scrivere una metaclasse. Una
\textbf{metaclasse} controlla la creazione dell'oggetto-classe stesso (e puo'
provocare conflitti di metaclasse nell'ereditarieta' multipla).
\code{\_\_subclasshook\_\_} interviene invece su \code{issubclass} e
\code{isinstance} e serve a implementare \emph{duck typing strutturale}: e' cosi'
che \code{collections.abc.Iterable} riconosce come sottoclasse qualunque tipo
definisca \dunder{iter}, senza ereditarieta' esplicita.
\end{soluzione}

\section{Soluzioni C}

\begin{soluzione}{D1 --- Dispatch dinamico scritto a mano}
\label{sol:D1}
\begin{lstlisting}[language=C]
typedef struct Shape Shape;

typedef struct {                     /* la vtable: una per CLASSE */
    void   (*draw)(const Shape *self);
    double (*area)(const Shape *self);
    void   (*dtor)(Shape *self);     /* distruttore virtuale */
    const char *type_name;           /* poor man's RTTI */
} ShapeVT;

struct Shape { const ShapeVT *vptr; };               /* "classe base" */
typedef struct { Shape base; double r; } Circle;     /* "classe derivata" */

static void   circle_draw(const Shape *s)
{ printf("Circle r=%g\n", ((const Circle *)s)->r); }

static double circle_area(const Shape *s)
{ double r = ((const Circle *)s)->r; return 3.14159265 * r * r; }

static void   circle_dtor(Shape *s) { free(s); }

static const ShapeVT CIRCLE_VT = { circle_draw, circle_area, circle_dtor, "Circle" };

Circle *circle_new(double r) {
    Circle *c = malloc(sizeof *c);
    c->base.vptr = &CIRCLE_VT;       /* <-- il "costruttore" installa il vptr */
    c->r = r;
    return c;
}

void draw_all(Shape *const *v, size_t n) {
    for (size_t i = 0; i < n; i++)
        v[i]->vptr->draw(v[i]);      /* == v[i]->draw() del C++ */
}
\end{lstlisting}

\textbf{Primo membro.} Lo standard garantisce che una struct non abbia padding
\emph{all'inizio}, quindi \code{(Shape *)circle == \&circle->base}: la
conversione derivata/base e' un no-op e il \code{this} non va aggiustato. E'
esattamente il caso dell'ereditarieta' singola in C++, quello senza thunk (A5).

\textbf{Il vptr lo inizializzi tu.} Se dimentichi quella riga hai un puntatore
indeterminato e la prima chiamata virtuale salta nel vuoto. In C++ e' il
compilatore a emettere l'assegnamento all'inizio di \emph{ogni} costruttore, e a
riscriverlo a ogni livello della gerarchia (che e' poi il motivo di A1).

\textbf{Distruttore virtuale} = una entry \code{dtor} nella vtable, invocata
sempre come \code{s->vptr->dtor(s)}; ogni ``derivata'' libera i propri membri e
poi chiama esplicitamente il distruttore della base.
\textbf{\code{dynamic\_cast}} = confronto del \code{vptr} (o del tag
\code{type\_name}) con quello atteso; per gerarchie a piu' livelli si mette
nella vtable un puntatore \code{parent} e si risale la catena.
\end{soluzione}

\begin{soluzione}{D2 --- Il cast che sembra innocuo}
\label{sol:D2}
\code{Circle *} e \code{Shape *} sono tipi di puntatore \textbf{non
compatibili}: la conversione \code{Circle **} $\rightarrow$ \code{Shape **} non
e' mai implicita e, una volta forzata col cast, l'accesso a \code{v[i]} legge un
oggetto il cui \emph{effective type} e' \code{Circle *} attraverso un lvalue di
tipo \code{Shape *}. E' una violazione della regola di \emph{aliasing}
(\emph{C11 6.5p7}), quindi undefined behavior --- innocuo nella pratica perche'
tutti i puntatori a struct hanno la stessa rappresentazione, ma il compilatore e'
autorizzato a riordinare o eliminare gli accessi.

La garanzia ``indirizzo del primo membro = indirizzo della struct''
(\emph{6.7.2.1p15}) riguarda \textbf{l'indirizzo dell'oggetto}, non la
compatibilita' fra \emph{tipi di puntatore}: dice che il cast
\code{Circle *} $\rightarrow$ \code{Shape *} da' l'indirizzo giusto, non che un
array di \code{Circle *} \emph{sia} un array di \code{Shape *}. La regola della
\emph{common initial sequence} (\emph{6.5.2.3}) e' poi enunciata solo per struct
che condividono una \code{union} visibile nella stessa unita' di traduzione.

Soluzione corretta: costruire un vero array di \code{Shape *}, convertendo
elemento per elemento (\code{arr[i] = \&c[i]->base;}). La conversione va fatta
\emph{sull'oggetto}, non sul puntatore a puntatore --- ed e' lo stesso motivo per
cui in Java gli array covarianti producono \code{ArrayStoreException} (E3).
\end{soluzione}

\begin{soluzione}{D3 --- Traduzione letterale di virtual}
\label{sol:D3}
\begin{itemize}
    \item \code{this} diventa il primo parametro esplicito \code{Shape *self};
    \item \code{const} diventa \code{const Shape *self}: nessuna differenza a
          runtime, solo type checking a compile time;
    \item la chiamata \code{s->area()} diventa \code{s->vptr->area(s)}: due
          dereferenziazioni (vptr, poi lo slot) piu' una chiamata indiretta;
    \item la \textbf{vtable e' un oggetto statico unico per classe}
          (\code{static const ShapeVT}), non per istanza. L'``ereditarieta''' della
          vtable e' letteralmente: copiare la tabella della base e sovrascrivere
          gli slot ridefiniti, \emph{mantenendo lo stesso indice} per lo stesso
          metodo. E' questa invarianza dell'indice a rendere il dispatch $O(1)$
          e indipendente dalla profondita' della gerarchia.
\end{itemize}
\textbf{Costo:} un puntatore per \emph{oggetto} (8 byte su architettura a 64
bit, piu' l'eventuale allineamento) e una tabella di $N$ puntatori per
\emph{classe}. Nessun altro overhead: e' il \emph{zero-overhead principle} ---
paghi solo se usi \code{virtual}, ed e' il motivo per cui in C++ il default e'
non virtuale (E2).
\end{soluzione}

\begin{soluzione}{D4 --- Due polimorfismi diversi}
\label{sol:D4}
\code{struct file\_operations} e' \textbf{overriding}: e' una vtable a tutti gli
effetti, condivisa da tutte le istanze di un ``tipo'' di file (ext4, procfs, un
driver di periferica), raggiunta da \code{struct file} attraverso il campo
\code{f\_op}. Ogni filesystem ridefinisce solo le entry che gli servono e lascia
\code{NULL} le altre (che il VFS interpreta come ``comportamento di default'':
un modo per emulare i metodi non astratti della base).

Il comparatore di \code{qsort} e' invece una \textbf{strategy}/callback: la
funzione e' passata per parametro alla singola chiamata e non e' associata al
\emph{tipo} degli oggetti ordinati; non esiste alcuna entita' ``classe'' che la
erediti. E' il polimorfismo della \emph{Strategy} del GoF, non quello del
\emph{Template Method}.

\code{container\_of(ptr, type, member)} risale dall'indirizzo di un membro
all'indirizzo della struct che lo contiene, sottraendo \code{offsetof}. E' il
\textbf{downcast} da base a derivata (\code{static\_cast<Derived *>(pBase)}) nel
caso in cui il sotto-oggetto base \emph{non} sia il primo membro: la stessa
aritmetica di offset che genererebbe il compilatore C++.
\end{soluzione}

\begin{soluzione}{D5 --- Perche' il C non puo' avere overriding automatico}
\label{sol:D5}
Perche' in C ogni chiamata di funzione e' \textbf{early binding}: il nome viene
risolto dal linker a \emph{un} singolo simbolo, e non esiste alcun canale --- ne'
un vptr implicito ne' metadati di tipo --- che leghi un \emph{oggetto} a un
insieme di funzioni. \code{static inline} negli header duplica la definizione in
ogni unita' di traduzione, ma resta comunque una risoluzione statica per nome.

Si aggiungono due impedimenti strutturali:
\begin{itemize}
    \item senza \emph{name mangling} non potresti nemmeno avere due funzioni
          \code{area} con lo stesso nome e parametri diversi nello stesso
          programma (niente overload, quindi a maggior ragione niente override);
    \item senza costruttori non esiste un punto garantito in cui installare il
          vptr: ogni schema di dispatch dinamico va costruito \emph{e mantenuto}
          a mano dal programmatore.
\end{itemize}
Ed e' precisamente quello che fanno i grandi progetti C: GObject/GTK, il VFS del
kernel Linux, COM su Windows. Il C non e' ``incapace'' di polimorfismo: si
limita a non automatizzarlo.
\end{soluzione}

\section{Soluzioni trasversali}

\begin{soluzione}{E1 --- Tabella comparativa del dispatch}
\label{sol:E1}
\begin{center}
\renewcommand{\arraystretch}{1.3}
\setlength{\tabcolsep}{4pt}
\scriptsize
\begin{tabular}{|>{\raggedright\arraybackslash}p{1.8cm}|>{\raggedright\arraybackslash}p{3.1cm}|>{\raggedright\arraybackslash}p{3.1cm}|>{\raggedright\arraybackslash}p{3.1cm}|>{\raggedright\arraybackslash}p{3.1cm}|}
\hline
& \textbf{C++} & \textbf{Java} & \textbf{Python} & \textbf{C} \\
\hline
Meccanismo &
vtable per classe + vptr per oggetto; indice noto a compile time &
\code{invokevirtual} (vtable) / \code{invokeinterface} (itable); JIT con inline cache e CHA &
lookup nel \code{\_\_dict\_\_} lungo l'MRO del tipo, con \emph{type method cache}, poi protocollo dei descrittori &
manuale: struct di puntatori a funzione \\
\hline
Costo &
2 load + chiamata indiretta; devirtualizzabile con \code{final}/LTO &
come C++, ma spesso \emph{inlinato per speculazione} e deottimizzato se serve &
il piu' alto: lookup hash + creazione del bound method (mitigato da cache e \code{LOAD\_METHOD}) &
identico al C++: e' la stessa cosa scritta a mano \\
\hline
Metodi a runtime &
no: layout fissato a compile time &
no per una classe gia' caricata (solo agent, redefinition, proxy dinamici) &
\textbf{si', sempre}: classi mutabili, \dunder{class} riassegnabile &
si', se la vtable non e' \code{const} \\
\hline
Method hiding &
si': name hiding degli overload, campi, membri non virtuali &
si': campi, \code{static}, \code{private} &
non esiste hiding, solo shadowing istanza/classe/MRO &
non applicabile \\
\hline
Chiamare la base &
\code{Base::f()} (chiamata non virtuale) &
\code{super.f()} $\rightarrow$ \code{invokespecial}; \code{B.super.f()} per i default &
\code{super().f()} (segue l'MRO!) oppure \code{A.f(self)} &
\code{parent\_vt->f(self)} esplicito \\
\hline
\end{tabular}
\end{center}
\end{soluzione}

\begin{soluzione}{E2 --- Virtuale per default oppure no}
\label{sol:E2}
Il C++ segue lo \textbf{zero-overhead principle}: chi non usa il polimorfismo
non deve pagare ne' il vptr (che oltretutto romperebbe la \emph{layout
compatibility} con il C e la natura POD dei tipi) ne' la chiamata indiretta. Per
questo \code{virtual} e' opt-in. Il prezzo e' che rendere virtuale un metodo in
un secondo momento e' un \textbf{cambio di ABI}, e che dimenticare
\code{virtual \~{}Base()} e' undefined behavior.

Java sceglie il default virtuale per uniformita' e perche' il costo e'
recuperabile \emph{a runtime}: la JVM esegue \textbf{Class Hierarchy Analysis}
(se una sola classe caricata implementa il metodo, la chiamata viene
devirtualizzata e inlinata, con una \emph{guard} e una deottimizzazione se in
seguito viene caricata una seconda implementazione) e usa \textbf{inline cache}
monomorfiche o bimorfiche. E' un'ottimizzazione che il C++ puo' fare solo con
LTO o con \code{final}, perche' la compilazione e' separata e l'insieme delle
classi non e' chiuso.

Il \textbf{fragile base class problem} taglia in entrambe le direzioni: con il
default virtuale ogni metodo pubblico diventa un punto di estensione non
intenzionale, e una modifica interna alla base puo' rompere sottoclassi che ne
fanno l'override (il caso classico e' \code{HashSet.addAll} che internamente
chiama \code{add}); con il default non virtuale il rischio si sposta
sull'\emph{hiding} silenzioso, che il compilatore non segnala. Python porta il
compromesso all'estremo opposto: massima flessibilita' (monkey patching, mixin
cooperativi, classi mutabili), nessuna garanzia statica e il costo di dispatch
piu' alto dei quattro.
\end{soluzione}

\begin{soluzione}{E3 --- Covarianza e controvarianza}
\label{sol:E3}
\textbf{In nessuno dei quattro linguaggi} \code{process(Derived)} sostituisce
\code{process(Base)} in modo corretto, ma per ragioni diverse.

In C++ e Java e' un \textbf{overload} con firma diversa: la derivata eredita
comunque \code{process(Base)} e vi aggiunge un secondo metodo (in Java
\code{@Override} darebbe errore; in C++ scatta anche il name hiding, A3). In
Python la firma non fa parte dell'identita' del metodo, quindi
\code{def process(self, d)} \emph{sostituisce} davvero quella della base: li'
l'override avviene, ma \textbf{viola LSP} e i type checker (mypy, pyright) lo
segnalano. In C il concetto non esiste.

Il principio di sostituzione di Liskov richiede parametri
\textbf{controvarianti} (la derivata deve accettare \emph{almeno} tutto cio' che
accetta la base) e return \textbf{covarianti} (deve promettere \emph{almeno}
quanto la base). Restringere il parametro a \code{Derived} e' covarianza dei
parametri, e non e' type-safe: con
\code{Base b = new Derived(); b.process(new Base());} il metodo riceverebbe un
oggetto che non sa trattare. E' lo stesso buco degli array covarianti di Java
(\code{ArrayStoreException}) e dei parametri covarianti di Eiffel, entrambi
riconosciuti come errori di progetto.

La controvarianza dei parametri sarebbe invece \emph{sicura}, ma quasi nessun
linguaggio la implementa perche' complica enormemente l'overload resolution e
nella pratica serve raramente: Python e TypeScript la accettano solo a livello
di type checker, Scala e Kotlin la esprimono sui generici (\code{-T}, \code{in T})
invece che sui parametri dei metodi.
\end{soluzione}

\begin{soluzione}{E4 --- Sette termini da non confondere}
\label{sol:E4}
\begin{itemize}
    \item \textbf{Overriding} --- ridefinizione, in una sottoclasse, di un
          metodo con la \emph{stessa firma}, risolta sul tipo \emph{dinamico}.
          C++ (solo se \code{virtual}), Java (default), Python (sempre). In C
          solo emulato a mano.
    \item \textbf{Overloading} --- piu' funzioni con lo stesso nome e parametri
          diversi, risolte a \emph{tempo di compilazione}. C++ e Java. In Python
          \textbf{no}: l'ultima definizione vince, si emula con
          \code{functools.singledispatch}. In C no.
    \item \textbf{Hiding / shadowing} --- un nome in uno scope piu' interno
          oscura quello esterno, \emph{senza} polimorfismo: campi e metodi
          \code{static} in Java, name hiding degli overload e membri non
          virtuali in C++. In Python esiste solo lo shadowing
          istanza/classe/MRO.
    \item \textbf{Dominanza} --- solo C++, con ereditarieta' virtuale: un membro
          di una classe piu' derivata domina quello della base virtuale e
          risolve un'ambiguita' che altrimenti sarebbe tale (A4).
    \item \textbf{Bridge method} --- solo Java: metodo sintetico generato dal
          compilatore per riconciliare l'erasure dei generici e i return
          covarianti con la firma che la JVM richiede per il dispatch (B5).
    \item \textbf{Thunk} --- C++ (e implementazioni C manuali): trampolino che
          aggiusta \code{this} nell'ereditarieta' multipla o virtuale prima di
          saltare all'override (A5).
    \item \textbf{Final overrider} --- C++: per ogni funzione virtuale e ogni
          classe, l'unica implementazione effettivamente installata in vtable;
          se non e' unica il programma e' ill-formed. Il concetto esiste anche
          in Java, ma e' sempre banale per via dell'ereditarieta' singola.
\end{itemize}
\end{soluzione}

%==============================================================
\chapter{Ripasso lampo}
%==============================================================

Le dieci affermazioni da avere in testa entrando all'esame.

\begin{enumerate}
    \item In C++ il tipo dinamico \emph{cresce} durante la costruzione e
          \emph{decresce} durante la distruzione: virtuali chiamate da
          ctor/dtor risolvono alla classe corrente. In Java il tipo dinamico e'
          quello finale fin dall'inizio, per questo la stessa situazione produce
          campi non ancora inizializzati.
    \item Overload $=$ statico, override $=$ dinamico. Se cambia la firma, non
          e' un override: e' un overload (e in C++ nasconde pure gli altri).
    \item Sono risolti sul \textbf{tipo statico}: campi, metodi \code{static},
          argomenti di default, controllo di accesso, scelta dell'overload,
          aritmetica dei puntatori negli array. Sul tipo dinamico c'e' solo il
          \emph{corpo} del metodo virtuale.
    \item Il controllo di accesso e' ortogonale al dispatch: un
          \code{private virtual} si puo' sovrascrivere ma non chiamare (NVI).
    \item Distruttore virtuale ogni volta che si distrugge tramite puntatore
          alla base; ma \code{delete[]} polimorfico resta UB comunque.
    \item Base non virtuale $=$ incorporata, una copia per ramo, offset
          costante. Base virtuale $=$ condivisa, una copia sola, offset letto a
          runtime: percio' la costruisce la classe \emph{most derived} e percio'
          \code{static\_cast} dalla base virtuale alla derivata e' vietato.
    \item In Java i generici sono cancellati: i \emph{bridge method} sono la
          toppa, e l'\emph{heap pollution} il suo effetto collaterale.
    \item Nel diamante dei default method vince la classe (\emph{class wins});
          fra interfacce non correlate non vince nessuno e si deve fare
          l'override esplicito con \code{X.super.f()}.
    \item In Python \code{super()} significa ``il prossimo nell'MRO
          dell'oggetto'', non ``la superclasse''; i dunder impliciti si cercano
          sul tipo, mai sull'istanza; il doppio underscore fa mangling e produce
          metodi di fatto non virtuali.
    \item In C il polimorfismo si scrive a mano: struct di puntatori a funzione
          come vtable, sotto-oggetto base come primo membro, vptr installato dal
          costruttore, \code{container\_of} al posto del downcast.
    \item Liskov: parametri controvarianti (o invarianti), return covarianti,
          precondizioni non piu' forti, postcondizioni non piu' deboli. Ogni
          override che viola questo schema e' un bug in attesa.
\end{enumerate}

\end{document}
